\appendixes

\begin{addendum}
	\chapter{Historias de usuario}
	\label{anx:hu}

		\subsubsection*{Épica 1: Gestión de Identidad y Acceso}

		\begin{userstory}[]
			\storyname{Autenticación y Autogestión de cuenta}
			\storyepic{Gestión de Identidad y Acceso}
			\storyrfs{1, 2, 3, 4, 5, 8}
			\storypriority{Alta}
			\storyrisk{Media}
			\storypoints{0.4}
			\storyprogrammer{Jose Luis Echemendia López}
			\storydescription{El usuario del sistema debe poder iniciar sesión, cerrar sesión, gestionar su perfil y recuperar su acceso para garantizar la seguridad de sus datos y su identidad en la plataforma.}
			\storyacceptancecriteria{
			\begin{itemize}
			\item Permitir registro e inicio con correo y contraseña.
			\item Procesar recuperación de acceso mediante \textit{token} enviado por correo.
			\item Permitir la edición del perfil propio sin intervención administrativa.
			\end{itemize}
			}
			\storyobservation{La historia concentra autenticación, recuperación de acceso y autogestión del perfil como capacidades base del módulo de identidad. Campos de entrada asociados: RF\#1 y RF\#3 (correo electrónico, contraseña), RF\#5 (token y nueva contraseña; o contraseña actual y nueva contraseña), RF\#8 (correo electrónico opcional).}
		\end{userstory}

		\begin{userstory}[]
			\storyname{Administración de Usuarios y Permisos}
			\storyepic{Gestión de Identidad y Acceso}
			\storyrfs{6, 7, 9, 10, 11}
			\storypriority{Alta}
			\storyrisk{Media}
			\storypoints{0.3}
			\storyprogrammer{Jose Luis Echemendia López}
			\storydescription{El administrador debe poder gestionar el listado de usuarios, sus roles y su estado activo o inactivo para controlar quién accede a la plataforma y qué acciones puede realizar según la jerarquía estática.}
			\storyacceptancecriteria{
			\begin{itemize}
			\item Listar y filtrar usuarios por rol o nombre.
			\item Modificar el rol asignado a cada usuario.
			\item Desactivar y activar cuentas sin eliminar la información persistida.
			\end{itemize}
			}
			\storyobservation{La gestión de cuentas debe preservar la trazabilidad administrativa y evitar borrados físicos sobre usuarios existentes. Campos de entrada asociados: RF\#9 (identificador de usuario en ruta y rol en consulta).}
		\end{userstory}

		\subsubsection*{Épica 2: Gestión de Dominios Curriculares}

		\begin{userstory}[]
			\storyname{Gestión de Dominios de Conocimiento}
			\storyepic{Gestión de Dominios Curriculares}
			\storyrfs{12, 13, 14, 15, 16, 17, 18, 19}
			\storypriority{Alta}
			\storyrisk{Media}
			\storypoints{0.4}
			\storyprogrammer{Jose Luis Echemendia López}
			\storydescription{El administrador debe poder crear, actualizar y supervisar los dominios curriculares para establecer el contexto de alto nivel donde se alojarán los \ac{FM}.}
			\storyacceptancecriteria{
			\begin{itemize}
			\item Crear dominios con nombre y descripción.
			\item Activar y desactivar dominios conforme a su estado de uso.
			\item Obtener cada dominio junto con sus modelos asociados.
			\end{itemize}
			}
			\storyobservation{El dominio funciona como unidad de organización para agrupar y consultar los modelos vinculados. Campos de entrada asociados: RF\#14 (nombre, descripción opcional) y RF\#15 (identificador de dominio en ruta; nombre y descripción opcionales).}
		\end{userstory}

		\subsubsection*{Épica 3: Ciclo de Vida del \ac{FM}}

		\begin{userstory}[]
			\storyname{Gestión Operativa de \ac{FM}}
			\storyepic{Ciclo de Vida del \ac{FM}}
			\storyrfs{20, 21, 22, 23, 24, 25}
			\storypriority{Alta}
			\storyrisk{Alta}
			\storypoints{1.0}
			\storyprogrammer{Jose Luis Echemendia López}
			\storydescription{El diseñador de modelo debe poder crear, listar y modificar los metadatos de los modelos de características dentro de un dominio para mantener un catálogo organizado de planes de estudio o arquitecturas curriculares.}
			\storyacceptancecriteria{
			\begin{itemize}
			\item Crear modelos asociados a un dominio válido.
			\item Actualizar metadatos preservando la integridad del modelo.
			\item Ofrecer paginación y filtrado para la consulta de modelos.
			\end{itemize}
			}
			\storyobservation{Esta historia cubre la administración operativa del catálogo de modelos, sin entrar todavía en su estructura interna. Campos de entrada asociados: RF\#21 (nombre, descripción opcional, identificador de dominio) y RF\#23 (identificador de modelo en ruta; nombre, descripción y estado activo opcionales).}
		\end{userstory}

		\begin{userstory}[]
			\storyname{Gestión de Versiones y Publicación}
			\storyepic{Ciclo de Vida del \ac{FM}}
			\storyrfs{26, 27, 28, 29, 30, 31, 32}
			\storypriority{Alta}
			\storyrisk{Alta}
			\storypoints{2.5}
			\storyprogrammer{Jose Luis Echemendia López}
			\storydescription{El diseñador de modelo y el revisador deben poder crear versiones, publicar instantáneas y archivar versiones antiguas para garantizar la trazabilidad y evolución del currículo sin afectar configuraciones previas.}
			\storyacceptancecriteria{
			\begin{itemize}
			\item Generar nuevas versiones mediante técnica de copia al escribir.
			\item Validar la versión antes de su publicación oficial.
			\item Permitir la recuperación de versiones archivadas.
			\end{itemize}
			}
			\storyobservation{La publicación debe conservar el historial del modelo y permitir volver a estados previos sin pérdida de información. Campos de entrada asociados: RF\#28 (identificador de modelo en ruta; identificador de versión origen y descripción opcionales).}
		\end{userstory}

		\begin{userstory}[]
			\storyname{Análisis de Métricas y Exportación}
			\storyepic{Ciclo de Vida del \ac{FM}}
			\storyrfs{33, 34}
			\storypriority{Media}
			\storyrisk{Media}
			\storypoints{0.8}
			\storyprogrammer{Jose Luis Echemendia López}
			\storydescription{El diseñador y el observador deben poder obtener estadísticas y exportar el modelo en múltiples formatos para compartirlo con otras herramientas o analizar su complejidad. Los formatos a soportar son:
				\begin{itemize}
					\item \ac{UVL}.
					\item \ac{XML}.
					\item \ac{TVL}.
					\item \ac{DIMACS}.
					\item \ac{JSON}.
					\item \ac{DOT}.
					\item mermaid.
				\end{itemize}
			}
			\storyacceptancecriteria{
			\begin{itemize}
			\item Calcular métricas de la versión publicada más reciente o de una versión específica.
			\item Exportar el modelo en formatos técnicos como \ac{UVL} y \textit{dimacs}.
			\item Exportar el modelo en formatos visuales como \textit{mermaid} y \textit{dot}.
			\item Guardar el resultado de la exportación en un archivo.
			\end{itemize}
			}
			\storyobservation{La exportación debe ser consistente con la versión consultada y servir tanto para intercambio técnico como para inspección visual.}
		\end{userstory}

		\subsubsection*{Épica 4: Modelado de Variabilidad Estructural}

		\begin{userstory}[]
			\storyname{Edición de la Jerarquía de Características}
			\storyepic{Modelado de Variabilidad Estructural}
			\storyrfs{53, 54, 55, 56, 60, 61, 62, 63}
			\storypriority{Alta}
			\storyrisk{Alta}
			\storypoints{1.0}
			\storyprogrammer{Jose Luis Echemendia López}
			\storydescription{El diseñador y el editor de modelo deben poder crear, mover, actualizar y desactivar características para construir el árbol de variabilidad que representa el plan de estudios.}
			\storyacceptancecriteria{
			\begin{itemize}
			\item Reubicar una \textit{feature} cambiando su padre sin perder consistencia.
			\item Recuperar subárboles completos de una característica.
			\item Aplicar desactivación lógica con soporte de versionado.
			\end{itemize}
			}
			\storyobservation{Las operaciones sobre la jerarquía deben preservar integridad estructural y trazabilidad de cambios. Campos de entrada asociados: RF\#53 (nombre, tipo, identificador de versión, y opcionales de padre, grupo, recurso y propiedades) y RF\#60 (identificador de característica en ruta; nombre, tipo, padre y grupo opcionales).}
		\end{userstory}

		\begin{userstory}[]
			\storyname{Modelado de Grupos y Relaciones Estructurales}
			\storyepic{Modelado de Variabilidad Estructural}
			\storyrfs{64, 65, 66, 67, 68, 69, 70, 71, 72, 73, 74, 75}
			\storypriority{Alta}
			\storyrisk{Alta}
			\storypoints{1.0}
			\storyprogrammer{Jose Luis Echemendia López}
			\storydescription{El diseñador y el editor de modelo deben poder definir grupos de características, como XOR y OR, y relaciones jerárquicas para restringir cómo se pueden combinar los módulos educativos.}
			\storyacceptancecriteria{
			\begin{itemize}
			\item Crear grupos con validaciones sobre el tipo de agrupación.
			\item Definir relaciones semánticas entre características.
			\item Permitir la definición de cardinalidades mínimas y máximas.
			\item Mantener consistencia al activar, desactivar o actualizar dichas estructuras.
			\end{itemize}
			}
			\storyobservation{Esta historia cubre la semántica estructural del árbol y la coherencia entre agrupaciones y relaciones. Campos de entrada asociados: RF\#66 (tipo de grupo, característica padre, cardinalidades), RF\#67 (identificador de grupo en ruta y campos opcionales), RF\#72 (tipo, origen, destino) y RF\#73 (identificador de relación en ruta y campos opcionales).}
		\end{userstory}

		\subsubsection*{Épica 5: Restricciones y Lógica Avanzada}

		\begin{userstory}[]
			\storyname{Gestión de Restricciones Transversales}
			\storyepic{Restricciones y Lógica Avanzada}
			\storyrfs{76, 77, 78, 79, 80, 81}
			\storypriority{Alta}
			\storyrisk{Alta}
			\storypoints{0.8}
			\storyprogrammer{Jose Luis Echemendia López}
			\storydescription{El diseñador de modelo debe poder crear y gestionar reglas de exclusión e implicación complejas para modelar dependencias que no son jerárquicas, por ejemplo, si se elige A, no se puede elegir B.}
			\storyacceptancecriteria{
			\begin{itemize}
			\item Permitir crear restricciones complejas con validación de sintaxis.
			\item Permitir expresiones lógicas complejas.
			\item Mantener control de versiones sobre las restricciones creadas.
			\end{itemize}
			}
			\storyobservation{Las restricciones deben integrarse con el modelo sin romper la evolución histórica del mismo. Campos de entrada asociados: RF\#78 (identificador de versión, expresión lógica y descripción opcional) y RF\#79 (identificador de restricción en ruta; expresión lógica y descripción opcionales).}
		\end{userstory}

		\begin{userstory}[]
			\storyname{Integración y Sincronización \ac{UVL}}
			\storyepic{Restricciones y Lógica Avanzada}
			\storyrfs{36, 37, 38, 39}
			\storypriority{Alta}
			\storyrisk{Alta}
			\storypoints{1.0}
			\storyprogrammer{Jose Luis Echemendia López}
			\storydescription{El diseñador deben poder operar directamente sobre el lenguaje \ac{UVL} o sincronizarlo desde la interfaz gráfica para mantener la potencia del modelado por código y la facilidad del modelado visual.}
			\storyacceptancecriteria{
			\begin{itemize}
			\item Sincronizar el modelo visual al editar el código \ac{UVL}.
			\item Persistir el \ac{UVL} analizado y validado.
			\item Aplicar \ac{UVL} a la estructura para regenerar el árbol desde el código \ac{UVL}.
			\end{itemize}
			}
			\storyobservation{La sincronización debe mantener equivalencia entre la representación textual y la representación visual del modelo. Campos de entrada asociados: RF\#36 y RF\#38 (identificador de modelo y versión en ruta, contenido \ac{UVL}).}
		\end{userstory}

		\subsubsection*{Épica 6: Validación y Análisis de Modelos}

		\begin{userstory}[]
			\storyname{Validación Lógica y Estructural del Modelo}
			\storyepic{Validación y Análisis de Modelos}
			\storyrfs{39, 40, 41, 42}
			\storypriority{Alta}
			\storyrisk{Alta}
			\storypoints{1.0}
			\storyprogrammer{Jose Luis Echemendia López}
			\storydescription{El revisador y el diseñador deben poder ejecutar análisis de consistencia, detectar características muertas y redundancias para asegurar que el plan de estudios sea lógicamente posible y no presente errores.}
			\storyacceptancecriteria{
			\begin{itemize}
			\item Evaluar la satisfacibilidad global del modelo mediante un solucionador \ac{SAT}.
			\item Detectar ciclos y nodos huérfanos.
			\item Identificar características muertas y redundancias.
			\end{itemize}
			}
			\storyobservation{El análisis debe advertir problemas estructurales y semánticos antes de publicar o derivar configuraciones.}
		\end{userstory}

		\subsubsection*{Épica 7: Configuración, Optimización y Recursos}

		\begin{userstory}[]
			\storyname{Gestión de Configuraciones}
			\storyepic{Configuración, Optimización y Recursos}
			\storyrfs{44, 45, 46, 47, 48, 49}
			\storypriority{Alta}
			\storyrisk{Alta}
			\storypoints{1.0}
			\storyprogrammer{Jose Luis Echemendia López}
			\storydescription{El configurador debe poder crear, actualizar y persistir selecciones de características para generar itinerarios educativos específicos para un perfil o estudiante.}
			\storyacceptancecriteria{
			\begin{itemize}
			\item Guardar configuraciones asociadas a una versión publicada.
			\item Ofrecer listado paginado de configuraciones almacenadas.
			\item Mantener operaciones de actualización y activación sobre configuraciones persistidas.
			\end{itemize}
			}
			\storyobservation{La configuración debe quedar ligada a una versión concreta del modelo para conservar su validez histórica. Campos de entrada asociados: RF\#44 (nombre, descripción opcional, identificador de versión y lista de características) y RF\#47 (identificador de configuración en ruta; nombre, descripción y lista de características opcionales).}
		\end{userstory}

		\begin{userstory}[]
			\storyname{Validación y Generación Automática de Configuraciones}
			\storyepic{Configuración, Optimización y Recursos}
			\storyrfs{49, 50, 51, 52}
			\storypriority{Alta}
			\storyrisk{Alta}
			\storypoints{1.0}
			\storyprogrammer{Jose Luis Echemendia López}
			\storydescription{El configurador y el estudiante deben poder validar si una selección es consistente y generar automáticamente opciones válidas u optimizadas para encontrar el mejor camino formativo sin violar prerrequisitos.}
			\storyacceptancecriteria{
			\begin{itemize}
			\item Validar selecciones sin necesidad de persistirlas.
			\item Ejecutar algoritmos de optimización para proponer mejores configuraciones.
			\end{itemize}
			}
			\storyobservation{La generación automática debe apoyar la toma de decisiones sin obligar a guardar resultados intermedios.}
		\end{userstory}

		\begin{userstory}[]
			\storyname{Enriquecimiento con Recursos y Taxonomías}
			\storyepic{Configuración, Optimización y Recursos}
			\storyrfs{57, 58, 59, 82, 83, 84, 85, 86, 87, 88}
			\storypriority{Media}
			\storyrisk{Media}
			\storypoints{0.8}
			\storyprogrammer{Jose Luis Echemendia López}
			\storydescription{El editor de modelo y el configurador deben poder etiquetar características y vincular recursos externos, como archivos o enlaces, para añadir contenido pedagógico real a la estructura lógica del modelo.}
			\storyacceptancecriteria{
			\begin{itemize}
			\item Asociar y desasociar etiquetas a características.
			\item Cargar recursos y asignar propietarios.
			\item Mantener versionada la información asociada a los recursos.
			\item Aceptar metadatos como idioma, licencia y otros de interés en el registro de los recursos.
			\end{itemize}
			}
			\storyobservation{Los recursos y etiquetas deben servir como apoyo semántico y documental para el uso del modelo. Campos de entrada asociados: RF\#84 (nombre, descripción opcional), RF\#87 (metadatos completos del recurso y contenido) y RF\#88 (identificador de recurso en ruta y campos opcionales, con restricción sobre propietario).}
		\end{userstory}


	\chapter{Tarjetas CRC}
	\label{anx:tCRC}

		\begin{crccard}[tb:crc_user]
			\crcclass{User}
			\crcresp
			{
				\begin{itemize}
				\item Representar usuarios y roles.
				\item Mantener estado activo/inactivo.
				\item Soportar auditoría (creado/actualizado/eliminado por).
				\end{itemize}
			}
			\crccolab
			{
				FeatureModel (owner/collaborator)\\
				FeatureModelCollaborator\\
				Repositorios de usuario
			}
		\end{crccard}

		\begin{crccard}[tb:crc_domain]
			\crcclass{Domain}
			\crcresp
			{
				\begin{itemize}
				\item Agrupar modelos por dominio.
				\item Mantener metadatos (nombre, descripción).
				\item Gestionar relación con modelos.
				\end{itemize}
			}
			\crccolab
			{
				FeatureModel
			}
		\end{crccard}

		\begin{crccard}[tb:crc_featuremodel]
			\crcclass{FeatureModel}
			\crcresp
			{
				\begin{itemize}
				\item Representar un \ac{FM} (nombre, descripción, dominio).
				\item Enlazar propietario y versiones.
				\item Mantener colaboradores.
				\end{itemize}
			}
			\crccolab
			{
				Domain\\
				User\\
				FeatureModelVersion\\
				FeatureModelCollaborator
			}
		\end{crccard}

		\begin{crccard}[tb:crc_featuremodelversion]
			\crcclass{FeatureModelVersion}
			\crcresp
			{
				\begin{itemize}
				\item Versionar un \ac{FM} (estado, snapshot, \ac{UVL}).
				\item Enlazar estructura completa (features, grupos, relaciones, restricciones y configuraciones).
				\end{itemize}
			}
			\crccolab
			{
				FeatureModel\\
				Feature\\
				FeatureGroup\\
				FeatureRelation\\
				Constraint\\
				Configuration
			}
		\end{crccard}

		\begin{crccard}[tb:crc_feature]
			\crcclass{Feature}
			\crcresp
			{
				\begin{itemize}
				\item Representar característica (nombre, tipo, propiedades).
				\item Mantener jerarquía padre/hijos.
				\item Enlazar grupo y recurso.
				\item Soportar tags y configuraciones.
				\end{itemize}
			}
			\crccolab
			{
				FeatureModelVersion\\
				FeatureGroup\\
				FeatureRelation\\
				Tag\\
				Configuration\\
				Resource
			}
		\end{crccard}

		\begin{crccard}[tb:crc_featuregroup]
			\crcclass{FeatureGroup}
			\crcresp
			{
				\begin{itemize}
				\item Definir agrupaciones XOR/OR.
				\item Controlar cardinalidades.
				\item Enlazar feature padre y miembros.
				\end{itemize}
			}
			\crccolab
			{
				FeatureModelVersion\\
				Feature
			}
		\end{crccard}

		\begin{crccard}[tb:crc_featurerelation]
			\crcclass{FeatureRelation}
			\crcresp
			{
				\begin{itemize}
				\item Representar relaciones cross-tree (requires/excludes).
				\end{itemize}
			}
			\crccolab
			{
				FeatureModelVersion\\
				Feature
			}
		\end{crccard}

		\begin{crccard}[tb:crc_constraint]
			\crcclass{Constraint}
			\crcresp
			{
				\begin{itemize}
				\item Persistir restricciones lógicas complejas (expr\_text, expr\_cnf).
				\end{itemize}
			}
			\crccolab
			{
				FeatureModelVersion
			}
		\end{crccard}

		\begin{crccard}[tb:crc_configuration]
			\crcclass{Configuration}
			\crcresp
			{
				\begin{itemize}
				\item Persistir selección de features.
				\item Enlazar tags.
				\item Mantener estado activo/inactivo.
				\end{itemize}
			}
			\crccolab
			{
				FeatureModelVersion\\
				Feature\\
				Tag
			}
		\end{crccard}

		\begin{crccard}[tb:crc_tag]
			\crcclass{Tag}
			\crcresp
			{
				\begin{itemize}
				\item Clasificar features/configuraciones.
				\item Mantener unicidad por nombre.
				\end{itemize}
			}
			\crccolab
			{
				Feature\\
				Configuration
			}
		\end{crccard}

		\begin{crccard}[tb:crc_resource]
			\crcclass{Resource}
			\crcresp
			{
				\begin{itemize}
				\item Representar recursos educativos.
				\item Mantener metadatos de contenido/licencia.
				\item Enlazar propietario.
				\end{itemize}
			}
			\crccolab
			{
				User\\
				Feature
			}
		\end{crccard}

		\begin{crccard}[tb:crc_fmversionmanager]
			\crcclass{FeatureModelVersionManager}
			\crcresp
			{
				\begin{itemize}
				\item Crear versiones (copy-on-write).
				\item Publicar/archivar/restaurar versiones.
				\item Generar snapshot y validaciones de transición.
				\end{itemize}
			}
			\crccolab
			{
				FeatureModelVersionRepository\\
				FeatureModelVersion\\
				Feature\\
				User\\
				FeatureModel
			}
		\end{crccard}

		\begin{crccard}[tb:crc_fmexportservice]
			\crcclass{FeatureModelExportService}
			\crcresp
			{
				\begin{itemize}
				\item Exportar \ac{FM} a \ac{UVL}/XML/TVL/DIMACS/\ac{JSON}/\ac{DOT}/Mermaid.
				\item Construir mapeo UUID $\leftrightarrow$ int.
				\end{itemize}
			}
			\crccolab
			{
				FeatureModelVersion\\
				Feature\\
				Constraint\\
				Enums de exportación
			}
		\end{crccard}

		\begin{crccard}[tb:crc_fmuvlimporter]
			\crcclass{FeatureModelUVLImporter}
			\crcresp
			{
				\begin{itemize}
				\item Validar \ac{UVL}.
				\item Construir estructura desde \ac{UVL}.
				\item Crear relaciones/restricciones y nueva versión.
				\end{itemize}
			}
			\crccolab
			{
				FeatureModelVersionManager\\
				Feature\\
				FeatureGroup\\
				FeatureRelation\\
				Constraint\\
				FeatureModelVersion\\
				User
			}
		\end{crccard}

		\begin{crccard}[tb:crc_fmlogicalvalidator]
			\crcclass{FeatureModelLogicalValidator}
			\crcresp
			{
				\begin{itemize}
				\item Ejecutar validación \ac{SAT}/\ac{SMT} (sympy/pysat/z3).
				\item Verificar consistencia global y configuraciones.
				\end{itemize}
			}
			\crccolab
			{
				InvalidConstraintException y otras excepciones\\
				Dependencias \ac{SAT}/\ac{SMT}
			}
		\end{crccard}

		\begin{crccard}[tb:crc_fmconfigurationgenerator]
			\crcclass{FeatureModelConfigurationGenerator}
			\crcresp
			{
				\begin{itemize}
				\item Generar configuraciones válidas.
				\item Aplicar estrategias GREEDY/RANDOM/BEAM/GENETIC/NSGA2.
				\end{itemize}
			}
			\crccolab
			{
				FeatureModelLogicalValidator\\
				GenerationStrategy\\
				Dependencias DEAP/OR-Tools/BDD
			}
		\end{crccard}

		\begin{crccard}[tb:crc_fmstructuralanalyzer]
			\crcclass{FeatureModelStructuralAnalyzer}
			\crcresp
			{
				\begin{itemize}
				\item Analizar estructura (dead features, redundancias, SCC).
				\item Calcular métricas de complejidad.
				\end{itemize}
			}
			\crccolab
			{
				AnalysisType\\
				networkx\\
				Excepciones estructurales
			}
		\end{crccard}

		\begin{crccard}[tb:crc_fmtreebuilder]
			\crcclass{FeatureModelTreeBuilder}
			\crcresp
			{
				\begin{itemize}
				\item Construir árbol completo para respuesta.
				\item Generar estadísticas.
				\item Generar \ac{UVL} efectivo.
				\end{itemize}
			}
			\crccolab
			{
				FeatureModelVersion\\
				Feature\\
				FeatureGroup\\
				FeatureModelExportService\\
				FeatureModelCompleteResponse
			}
		\end{crccard}

		\begin{crccard}[tb:crc_fmanalysisfacade]
			\crcclass{FeatureModelAnalysisFacade}
			\crcresp
			{
				\begin{itemize}
				\item Orquestar análisis lógico y estructural.
				\item Calcular commonality/core/atomic sets.
				\item Ejecutar validación \ac{UVL} opcional (Flamapy).
				\end{itemize}
			}
			\crccolab
			{
				FeatureModelLogicalValidator\\
				FeatureModelStructuralAnalyzer\\
				FeatureModelUVLImporter
			}
		\end{crccard}

		\begin{crccard}[tb:crc_settingsservice]
			\crcclass{SettingsService}
			\crcresp
			{
				\begin{itemize}
				\item Consultar y actualizar settings.
				\item Cachear en Redis.
				\end{itemize}
			}
			\crccolab
			{
				AppSetting\\
				redis\_client
			}
		\end{crccard}

		\begin{crccard}[tb:crc_appsetting]
			\crcclass{AppSetting}
			\crcresp
			{
				\begin{itemize}
				\item Persistir claves de configuración.
				\item Mantener descripción y valor.
				\end{itemize}
			}de
			\crccolab
			{
				SettingsService\\
				Redis cache
			}
		\end{crccard}

	\chapter{Tareas de ingenierías}
	\label{anx:ti}

		\subsubsection*{HU-01: Autenticación y Autogestión de cuenta}

		\begin{engineeringtask}[tb:engtask_t001]
		\engtaskuserstory{1}
		\engtaskname{Diseñar modelo de token de recuperación}
		\engtasktype{Definir modelos}
		\engtaskpointestimation{6}
		\engtaskstartdate{20}{4}{2026}
		\engtaskenddate{20}{4}{2026}
		\engtaskdescription{Diseñar y documentar el modelo de token de recuperación de contraseña, incluyendo ciclo de vida, expiración, estado, relación con usuario y campos de auditoría compatibles con la estrategia de persistencia.}
		\engtaskprogrammer{Jose Luis Echemendia López}
		\end{engineeringtask}

		\begin{engineeringtask}[tb:engtask_t002]
		\engtaskuserstory{1}
		\engtaskname{Implementar esquemas de autenticación}
		\engtasktype{Crear esquemas de validación entrada de la api y salida}
		\engtaskpointestimation{8}
		\engtaskstartdate{20}{4}{2026}
		\engtaskenddate{21}{4}{2026}
		\engtaskdescription{Definir esquemas de entrada y salida para registro, inicio de sesión, recuperación de acceso, restablecimiento de contraseña y actualización de perfil, con validaciones explícitas y mensajes de error consistentes.}
		\engtaskprogrammer{Jose Luis Echemendia López}
		\end{engineeringtask}

		\begin{engineeringtask}[tb:engtask_t003]
		\engtaskuserstory{1}
		\engtaskname{Implementar servicio de autenticación y cuenta}
		\engtasktype{Crear servicio}
		\engtaskpointestimation{12}
		\engtaskstartdate{21}{4}{2026}
		\engtaskenddate{22}{4}{2026}
		\engtaskdescription{Diseñar e implementar el servicio de identidad para registro, autenticación, gestión de credenciales, emisión y verificación de tokens, y actualización de perfil, garantizando trazabilidad y cumplimiento de políticas de seguridad.}
		\engtaskprogrammer{Jose Luis Echemendia López}
		\end{engineeringtask}

		\begin{engineeringtask}[tb:engtask_t004]
		\engtaskuserstory{1}
		\engtaskname{Exponer endpoints de auth y perfil}
		\engtasktype{Crear controlador}
		\engtaskpointestimation{10}
		\engtaskstartdate{22}{4}{2026}
		\engtaskenddate{23}{4}{2026}
		\engtaskdescription{Implementar el controlador de identidad con endpoints de registro, inicio/cierre de sesión, solicitud de recuperación, restablecimiento de contraseña y edición de perfil propio; incorporar validación de solicitudes, manejo uniforme de errores y contratos de respuesta consistentes.}
		\engtaskprogrammer{Jose Luis Echemendia López}
		\end{engineeringtask}

		\begin{engineeringtask}[tb:engtask_t005]
		\engtaskuserstory{1}
		\engtaskname{Enviar correos de recuperación asíncronamente}
		\engtasktype{Crear tareas de celery}
		\engtaskpointestimation{6}
		\engtaskstartdate{23}{4}{2026}
		\engtaskenddate{24}{4}{2026}
		\engtaskdescription{Implementar la tarea asíncrona de notificación para recuperación de cuenta con políticas de reintento, plantillas de correo versionadas y registro de trazas operativas para auditoría y soporte.}
		\engtaskprogrammer{Jose Luis Echemendia López}
		\end{engineeringtask}

		\subsubsection*{HU-02: Administración de Usuarios y Permisos}

		\begin{engineeringtask}[tb:engtask_t006]
		\engtaskuserstory{2}
		\engtaskname{Implementar esquemas para gestión de usuarios}
		\engtasktype{Crear esquemas de validación entrada de la api y salida}
		\engtaskpointestimation{6}
		\engtaskstartdate{27}{4}{2026}
		\engtaskenddate{27}{4}{2026}
		\engtaskdescription{Definir contratos de entrada y salida para listado paginado, filtrado por rol y término de búsqueda, actualización de rol y cambio de estado de cuenta, con reglas de validación y paginación estandarizadas.}
		\engtaskprogrammer{Jose Luis Echemendia López}
		\end{engineeringtask}

		\begin{engineeringtask}[tb:engtask_t007]
		\engtaskuserstory{2}
		\engtaskname{Implementar repositorio administrativo de usuarios}
		\engtasktype{Crear repositorio}
		\engtaskpointestimation{8}
		\engtaskstartdate{27}{4}{2026}
		\engtaskenddate{28}{4}{2026}
		\engtaskdescription{Implementar el repositorio de usuarios administrativos con consultas paginadas, filtros combinables, búsqueda textual y operaciones de actualización de rol y estado, preservando eliminación lógica y consistencia transaccional.}
		\engtaskprogrammer{Jose Luis Echemendia López}
		\end{engineeringtask}

		\begin{engineeringtask}[tb:engtask_t008]
		\engtaskuserstory{2}
		\engtaskname{Implementar servicio de administración de usuarios}
		\engtasktype{Crear servicio}
		\engtaskpointestimation{10}
		\engtaskstartdate{28}{4}{2026}
		\engtaskenddate{29}{4}{2026}
		\engtaskdescription{Implementar el servicio de administración de usuarios para orquestar permisos, reglas de transición de rol y validaciones de seguridad, incorporando controles de autorización y registro de acciones críticas.}
		\engtaskprogrammer{Jose Luis Echemendia López}
		\end{engineeringtask}

		\begin{engineeringtask}[tb:engtask_t009]
		\engtaskuserstory{2}
		\engtaskname{Exponer endpoints de administración de usuarios}
		\engtasktype{Crear controlador}
		\engtaskpointestimation{8}
		\engtaskstartdate{29}{4}{2026}
		\engtaskenddate{30}{4}{2026}
		\engtaskdescription{Implementar el controlador administrativo para listado paginado, filtrado y búsqueda de usuarios, cambio de rol y activación/desactivación de cuentas; aplicar autorización por rol, validación de parámetros de ruta/consulta y trazabilidad de acciones administrativas.}
		\engtaskprogrammer{Jose Luis Echemendia López}
		\end{engineeringtask}

		\begin{engineeringtask}[tb:engtask_t010]
		\engtaskuserstory{2}
		\engtaskname{Definir excepciones de administración de usuarios}
		\engtasktype{Crear excepciones personalizadas}
		\engtaskpointestimation{4}
		\engtaskstartdate{30}{4}{2026}
		\engtaskenddate{30}{4}{2026}
		\engtaskdescription{Diseñar excepciones de dominio para transición inválida de rol, usuario inexistente y operación no autorizada, estandarizando códigos de error y semántica de respuesta para la capa API.}
		\engtaskprogrammer{Jose Luis Echemendia López}
		\end{engineeringtask}

		\subsubsection*{HU-03: Gestión de Dominios de Conocimiento}

		\begin{engineeringtask}[tb:engtask_t011]
		\engtaskuserstory{3}
		\engtaskname{Ajustar modelo de dominio curricular}
		\engtasktype{Definir modelos}
		\engtaskpointestimation{6}
		\engtaskstartdate{4}{5}{2026}
		\engtaskenddate{4}{5}{2026}
		\engtaskdescription{Refinar el modelo de dominio curricular, incorporando restricciones de unicidad, atributos de estado, campos de auditoría y relaciones con modelos para asegurar consistencia semántica y persistencia confiable.}
		\engtaskprogrammer{Jose Luis Echemendia López}
		\end{engineeringtask}

		\begin{engineeringtask}[tb:engtask_t012]
		\engtaskuserstory{3}
		\engtaskname{Implementar esquemas de dominio}
		\engtasktype{Crear esquemas de validación entrada de la api y salida}
		\engtaskpointestimation{6}
		\engtaskstartdate{4}{5}{2026}
		\engtaskenddate{5}{5}{2026}
		\engtaskdescription{Definir esquemas para creación, actualización, consulta, búsqueda y detalle de dominios con modelos asociados, estableciendo validaciones de formato, límites y contratos de serialización coherentes.}
		\engtaskprogrammer{Jose Luis Echemendia López}
		\end{engineeringtask}

		\begin{engineeringtask}[tb:engtask_t013]
		\engtaskuserstory{3}
		\engtaskname{Implementar repositorio de dominios}
		\engtasktype{Crear repositorio}
		\engtaskpointestimation{8}
		\engtaskstartdate{5}{5}{2026}
		\engtaskenddate{6}{5}{2026}
		\engtaskdescription{Implementar el repositorio de dominios con consultas paginadas, criterios de búsqueda por nombre y descripción, y recuperación eficiente de relaciones con modelos asociados.}
		\engtaskprogrammer{Jose Luis Echemendia López}
		\end{engineeringtask}

		\begin{engineeringtask}[tb:engtask_t014]
		\engtaskuserstory{3}
		\engtaskname{Implementar servicio de gobernanza de dominios}
		\engtasktype{Crear servicio}
		\engtaskpointestimation{8}
		\engtaskstartdate{6}{5}{2026}
		\engtaskenddate{7}{5}{2026}
		\engtaskdescription{Implementar el servicio de gobernanza de dominios para gestionar ciclo de vida, validaciones de unicidad e integridad referencial, y reglas de activación/desactivación con trazabilidad administrativa.}
		\engtaskprogrammer{Jose Luis Echemendia López}
		\end{engineeringtask}

		\begin{engineeringtask}[tb:engtask_t015]
		\engtaskuserstory{3}
		\engtaskname{Exponer endpoints de dominios}
		\engtasktype{Crear controlador}
		\engtaskpointestimation{8}
		\engtaskstartdate{7}{5}{2026}
		\engtaskenddate{8}{5}{2026}
		\engtaskdescription{Implementar el controlador de dominios para operaciones de creación, consulta, actualización, búsqueda y cambio de estado, incluyendo endpoint de detalle con modelos asociados; asegurar validación de datos, paginación y estandarización de respuestas HTTP.}
		\engtaskprogrammer{Jose Luis Echemendia López}
		\end{engineeringtask}

		\subsubsection*{HU-04: Gestión Operativa de FM}

		\begin{engineeringtask}[tb:engtask_t016]
		\engtaskuserstory{4}
		\engtaskname{Ajustar modelo de FeatureModel y colaboradores}
		\engtasktype{Definir modelos}
		\engtaskpointestimation{8}
		\engtaskstartdate{11}{5}{2026}
		\engtaskenddate{11}{5}{2026}
		\engtaskdescription{Consolidar el modelo de FeatureModel y colaboradores, formalizando relaciones con dominio y propietario, además de reglas de unicidad y consistencia para operaciones de lectura y escritura.}
		\engtaskprogrammer{Jose Luis Echemendia López}
		\end{engineeringtask}

		\begin{engineeringtask}[tb:engtask_t017]
		\engtaskuserstory{4}
		\engtaskname{Implementar esquemas operativos de FM}
		\engtasktype{Crear esquemas de validación entrada de la api y salida}
		\engtaskpointestimation{8}
		\engtaskstartdate{11}{5}{2026}
		\engtaskenddate{12}{5}{2026}
		\engtaskdescription{Definir esquemas operativos para alta, actualización, listado paginado y consulta de detalle del modelo, incluyendo metadatos obligatorios, validaciones de consistencia y formato de respuesta uniforme.}
		\engtaskprogrammer{Jose Luis Echemendia López}
		\end{engineeringtask}

		\begin{engineeringtask}[tb:engtask_t018]
		\engtaskuserstory{4}
		\engtaskname{Implementar repositorio de FM}
		\engtasktype{Crear repositorio}
		\engtaskpointestimation{8}
		\engtaskstartdate{12}{5}{2026}
		\engtaskenddate{13}{5}{2026}
		\engtaskdescription{Implementar el repositorio de modelos de características con operaciones de persistencia y consultas por dominio, estado y criterios de búsqueda, optimizando lectura para escenarios paginados.}
		\engtaskprogrammer{Jose Luis Echemendia López}
		\end{engineeringtask}

		\begin{engineeringtask}[tb:engtask_t019]
		\engtaskuserstory{4}
		\engtaskname{Implementar servicio operativo de FM}
		\engtasktype{Crear servicio}
		\engtaskpointestimation{10}
		\engtaskstartdate{13}{5}{2026}
		\engtaskenddate{14}{5}{2026}
		\engtaskdescription{Implementar el servicio operativo del catálogo de modelos para gestionar altas, cambios de metadatos y estado operativo, aplicando validaciones de negocio e integridad de información.}
		\engtaskprogrammer{Jose Luis Echemendia López}
		\end{engineeringtask}

		\begin{engineeringtask}[tb:engtask_t020]
		\engtaskuserstory{4}
		\engtaskname{Exponer endpoints de ciclo operativo de FM}
		\engtasktype{Crear controlador}
		\engtaskpointestimation{8}
		\engtaskstartdate{14}{5}{2026}
		\engtaskenddate{15}{5}{2026}
		\engtaskdescription{Implementar el controlador de modelos de características para alta, consulta, actualización y activación/desactivación, con soporte de filtros por dominio, paginación y políticas de autorización según rol funcional.}
		\engtaskprogrammer{Jose Luis Echemendia López}
		\end{engineeringtask}

		\subsubsection*{HU-05: Gestión de Versiones y Publicación}

		\begin{engineeringtask}[tb:engtask_t021]
		\engtaskuserstory{5}
		\engtaskname{Ajustar modelo de versionado de FM}
		\engtasktype{Definir modelos}
		\engtaskpointestimation{8}
		\engtaskstartdate{18}{5}{2026}
		\engtaskenddate{18}{5}{2026}
		\engtaskdescription{Definir el modelo de versionado con estados explícitos, metadatos de snapshot y trazabilidad completa, estableciendo reglas de transición válidas para sostener el ciclo de vida de versiones.}
		\engtaskprogrammer{Jose Luis Echemendia López}
		\end{engineeringtask}

		\begin{engineeringtask}[tb:engtask_t022]
		\engtaskuserstory{5}
		\engtaskname{Implementar repositorio de versiones}
		\engtasktype{Crear repositorio}
		\engtaskpointestimation{10}
		\engtaskstartdate{18}{5}{2026}
		\engtaskenddate{19}{5}{2026}
		\engtaskdescription{Implementar el repositorio de versiones con consultas por modelo, localización de versión publicada vigente y recuperación de versiones específicas con sus dependencias necesarias.}
		\engtaskprogrammer{Jose Luis Echemendia López}
		\end{engineeringtask}

		\begin{engineeringtask}[tb:engtask_t023]
		\engtaskuserstory{5}
		\engtaskname{Implementar gestor de versionado copy-on-write}
		\engtasktype{Crear clase}
		\engtaskpointestimation{14}
		\engtaskstartdate{19}{5}{2026}
		\engtaskenddate{21}{5}{2026}
		\engtaskdescription{Diseñar e implementar la clase FeatureModelVersionManager para orquestar creación de versiones por copy-on-write, publicación, archivado y restauración, con validación estricta de transiciones de estado.}
		\engtaskprogrammer{Jose Luis Echemendia López}
		\end{engineeringtask}

		\begin{engineeringtask}[tb:engtask_t024]
		\engtaskuserstory{5}
		\engtaskname{Exponer endpoints de versiones y publicación}
		\engtasktype{Crear controlador}
		\engtaskpointestimation{10}
		\engtaskstartdate{21}{5}{2026}
		\engtaskenddate{22}{5}{2026}
		\engtaskdescription{Implementar el controlador de versionado para listar/consultar versiones, crear versiones por copy-on-write, publicar, archivar, restaurar y obtener estructura completa de una versión; incluir validación de transiciones de estado y control de permisos por perfil.}
		\engtaskprogrammer{Jose Luis Echemendia López}
		\end{engineeringtask}

		\begin{engineeringtask}[tb:engtask_t025]
		\engtaskuserstory{5}
		\engtaskname{Definir excepciones de transición de versiones}
		\engtasktype{Crear excepciones personalizadas}
		\engtaskpointestimation{4}
		\engtaskstartdate{22}{5}{2026}
		\engtaskenddate{22}{5}{2026}
		\engtaskdescription{Definir excepciones de dominio para transiciones inválidas, condiciones de publicación no satisfechas y conflictos de estado, con mapeo consistente a respuestas HTTP y contexto de diagnóstico.}
		\engtaskprogrammer{Jose Luis Echemendia López}
		\end{engineeringtask}

		\subsubsection*{HU-06: Análisis de Métricas y Exportación}

		\begin{engineeringtask}[tb:engtask_t026]
		\engtaskuserstory{6}
		\engtaskname{Implementar servicio de exportación multiformato}
		\engtasktype{Crear servicio}
		\engtaskpointestimation{16}
		\engtaskstartdate{25}{5}{2026}
		\engtaskenddate{26}{5}{2026}
		\engtaskdescription{Implementar el servicio de exportación multiformato para generar representaciones UVL, XML, TVL, DIMACS, JSON, \ac{DOT} y Mermaid, incorporando mapeo estable de identificadores y validación de consistencia del artefacto generado.}
		\engtaskprogrammer{Jose Luis Echemendia López}
		\end{engineeringtask}

		\begin{engineeringtask}[tb:engtask_t027]
		\engtaskuserstory{6}
		\engtaskname{Implementar esquemas de métricas y exportación}
		\engtasktype{Crear esquemas de validación entrada de la api y salida}
		\engtaskpointestimation{6}
		\engtaskstartdate{26}{5}{2026}
		\engtaskenddate{26}{5}{2026}
		\engtaskdescription{Definir esquemas de solicitud y respuesta para consulta de métricas y exportación por versión objetivo o versión vigente publicada, incluyendo parámetros de formato y validaciones de compatibilidad.}
		\engtaskprogrammer{Jose Luis Echemendia López}
		\end{engineeringtask}

		\begin{engineeringtask}[tb:engtask_t028]
		\engtaskuserstory{6}
		\engtaskname{Exponer endpoints de estadísticas y exportación}
		\engtasktype{Crear controlador}
		\engtaskpointestimation{8}
		\engtaskstartdate{26}{5}{2026}
		\engtaskenddate{27}{5}{2026}
		\engtaskdescription{Implementar el controlador de análisis y exportación para consultar métricas por versión y generar artefactos en formatos técnicos y visuales; soportar negociación de formato, descarga/streaming seguro y manejo de errores de serialización.}
		\engtaskprogrammer{Jose Luis Echemendia López}
		\end{engineeringtask}

		\begin{engineeringtask}[tb:engtask_t029]
		\engtaskuserstory{6}
		\engtaskname{Implementar tareas asíncronas de exportación}
		\engtasktype{Crear tareas de celery}
		\engtaskpointestimation{8}
		\engtaskstartdate{27}{5}{2026}
		\engtaskenddate{28}{5}{2026}
		\engtaskdescription{Implementar tareas asíncronas de exportación para procesar cargas intensivas fuera del hilo de petición, con persistencia de estado, almacenamiento de artefactos y mecanismos de recuperación ante fallos.}
		\engtaskprogrammer{Jose Luis Echemendia López}
		\end{engineeringtask}

		\begin{engineeringtask}[tb:engtask_t030]
		\engtaskuserstory{6}
		\engtaskname{Definir excepciones de exportación}
		\engtasktype{Crear excepciones personalizadas}
		\engtaskpointestimation{4}
		\engtaskstartdate{28}{5}{2026}
		\engtaskenddate{28}{5}{2026}
		\engtaskdescription{Definir excepciones específicas para formato no soportado, versión inexistente y errores de serialización o generación de artefactos, con semántica clara para observabilidad y soporte operativo.}
		\engtaskprogrammer{Jose Luis Echemendia López}
		\end{engineeringtask}

		\subsubsection*{HU-07: Edición de la Jerarquía de Características}

		\begin{engineeringtask}[tb:engtask_t031]
		\engtaskuserstory{7}
		\engtaskname{Ajustar modelo de Feature para edición jerárquica}
		\engtasktype{Definir modelos}
		\engtaskpointestimation{8}
		\engtaskstartdate{1}{6}{2026}
		\engtaskenddate{1}{6}{2026}
		\engtaskdescription{Ajustar el modelo de Feature para soportar edición jerárquica segura, contemplando relaciones padre-hijo, pertenencia a grupo, vínculo con recursos, propiedades extensibles y estado lógico activo.}
		\engtaskprogrammer{Jose Luis Echemendia López}
		\end{engineeringtask}

		\begin{engineeringtask}[tb:engtask_t032]
		\engtaskuserstory{7}
		\engtaskname{Implementar esquemas de jerarquía de features}
		\engtasktype{Crear esquemas de validación entrada de la api y salida}
		\engtaskpointestimation{8}
		\engtaskstartdate{1}{6}{2026}
		\engtaskenddate{2}{6}{2026}
		\engtaskdescription{Definir esquemas de entrada y salida para creación, actualización, movimiento, consulta de detalle y subárbol, y cambio de estado de características, con validaciones estructurales y referenciales.}
		\engtaskprogrammer{Jose Luis Echemendia López}
		\end{engineeringtask}

		\begin{engineeringtask}[tb:engtask_t033]
		\engtaskuserstory{7}
		\engtaskname{Implementar repositorio de features}
		\engtasktype{Crear repositorio}
		\engtaskpointestimation{10}
		\engtaskstartdate{2}{6}{2026}
		\engtaskenddate{3}{6}{2026}
		\engtaskdescription{Implementar el repositorio de características con operaciones de recorrido jerárquico, consultas por versión, actualizaciones parciales y desactivación lógica, garantizando integridad de relaciones padre-hijo.}
		\engtaskprogrammer{Jose Luis Echemendia López}
		\end{engineeringtask}

		\begin{engineeringtask}[tb:engtask_t034]
		\engtaskuserstory{7}
		\engtaskname{Implementar servicio de jerarquía de features}
		\engtasktype{Crear servicio}
		\engtaskpointestimation{12}
		\engtaskstartdate{3}{6}{2026}
		\engtaskenddate{4}{6}{2026}
		\engtaskdescription{Implementar el servicio de jerarquía de características para validar movimientos en el árbol, prevenir inconsistencias estructurales y ejecutar versionado copy-on-write en operaciones de activación y desactivación.}
		\engtaskprogrammer{Jose Luis Echemendia López}
		\end{engineeringtask}

		\begin{engineeringtask}[tb:engtask_t035]
		\engtaskuserstory{7}
		\engtaskname{Exponer endpoints de features estructurales}
		\engtasktype{Crear controlador}
		\engtaskpointestimation{8}
		\engtaskstartdate{4}{6}{2026}
		\engtaskenddate{5}{6}{2026}
		\engtaskdescription{Implementar el controlador de características para creación, consulta de detalle/subárbol, actualización, movimiento en jerarquía y activación/desactivación lógica; reforzar validaciones de integridad estructural y autorización por rol de edición.}
		\engtaskprogrammer{Jose Luis Echemendia López}
		\end{engineeringtask}

		\subsubsection*{HU-08: Modelado de Grupos y Relaciones Estructurales}

		\begin{engineeringtask}[tb:engtask_t036]
		\engtaskuserstory{8}
		\engtaskname{Ajustar modelos de FeatureGroup y FeatureRelation}
		\engtasktype{Definir modelos}
		\engtaskpointestimation{8}
		\engtaskstartdate{8}{6}{2026}
		\engtaskenddate{8}{6}{2026}
		\engtaskdescription{Definir y refinar los modelos de grupos y relaciones, incorporando tipo semántico, cardinalidades, extremos de vínculo y estado lógico, junto con reglas de integridad para persistencia y validación.}
		\engtaskprogrammer{Jose Luis Echemendia López}
		\end{engineeringtask}

		\begin{engineeringtask}[tb:engtask_t037]
		\engtaskuserstory{8}
		\engtaskname{Implementar esquemas de grupos y relaciones}
		\engtasktype{Crear esquemas de validación entrada de la api y salida}
		\engtaskpointestimation{8}
		\engtaskstartdate{8}{6}{2026}
		\engtaskenddate{9}{6}{2026}
		\engtaskdescription{Definir contratos de API para listado, creación, actualización y cambio de estado de grupos y relaciones, incorporando validaciones de cardinalidad, tipos permitidos y consistencia entre nodos.}
		\engtaskprogrammer{Jose Luis Echemendia López}
		\end{engineeringtask}

		\begin{engineeringtask}[tb:engtask_t038]
		\engtaskuserstory{8}
		\engtaskname{Implementar repositorios de grupos y relaciones}
		\engtasktype{Crear repositorio}
		\engtaskpointestimation{10}
		\engtaskstartdate{9}{6}{2026}
		\engtaskenddate{10}{6}{2026}
		\engtaskdescription{Implementar repositorios de grupos y relaciones con consultas por versión, característica padre y tipo semántico, optimizando la carga de datos para validaciones estructurales posteriores.}
		\engtaskprogrammer{Jose Luis Echemendia López}
		\end{engineeringtask}

		\begin{engineeringtask}[tb:engtask_t039]
		\engtaskuserstory{8}
		\engtaskname{Implementar servicio de semántica estructural}
		\engtasktype{Crear servicio}
		\engtaskpointestimation{12}
		\engtaskstartdate{10}{6}{2026}
		\engtaskenddate{11}{6}{2026}
		\engtaskdescription{Implementar el servicio de semántica estructural para validar cardinalidades, coherencia de relaciones cross-tree y reglas de versionado al aplicar cambios sobre la estructura del modelo.}
		\engtaskprogrammer{Jose Luis Echemendia López}
		\end{engineeringtask}

		\begin{engineeringtask}[tb:engtask_t040]
		\engtaskuserstory{8}
		\engtaskname{Exponer endpoints de grupos y relaciones}
		\engtasktype{Crear controlador}
		\engtaskpointestimation{8}
		\engtaskstartdate{11}{6}{2026}
		\engtaskenddate{12}{6}{2026}
		\engtaskdescription{Implementar el controlador de grupos y relaciones para gestión completa de agrupaciones (incluyendo cardinalidades) y relaciones cross-tree; validar reglas semánticas del modelo y aplicar autorización granular para operaciones de edición.}
		\engtaskprogrammer{Jose Luis Echemendia López}
		\end{engineeringtask}

		\subsubsection*{HU-09: Gestión de Restricciones Transversales}

		\begin{engineeringtask}[tb:engtask_t041]
		\engtaskuserstory{9}
		\engtaskname{Ajustar modelo de restricciones lógicas}
		\engtasktype{Definir modelos}
		\engtaskpointestimation{6}
		\engtaskstartdate{15}{6}{2026}
		\engtaskenddate{15}{6}{2026}
		\engtaskdescription{Definir el modelo de restricciones lógicas con representación textual y normalizada, estado de vigencia y metadatos de versión, habilitando trazabilidad y análisis posterior.}
		\engtaskprogrammer{Jose Luis Echemendia López}
		\end{engineeringtask}

		\begin{engineeringtask}[tb:engtask_t042]
		\engtaskuserstory{9}
		\engtaskname{Implementar esquemas de restricciones}
		\engtasktype{Crear esquemas de validación entrada de la api y salida}
		\engtaskpointestimation{6}
		\engtaskstartdate{15}{6}{2026}
		\engtaskenddate{16}{6}{2026}
		\engtaskdescription{Definir esquemas de solicitud y respuesta para ciclo completo de restricciones, incluyendo validaciones sintácticas de expresión, serialización uniforme y estructura de errores de dominio.}
		\engtaskprogrammer{Jose Luis Echemendia López}
		\end{engineeringtask}

		\begin{engineeringtask}[tb:engtask_t043]
		\engtaskuserstory{9}
		\engtaskname{Implementar repositorio de restricciones}
		\engtasktype{Crear repositorio}
		\engtaskpointestimation{8}
		\engtaskstartdate{16}{6}{2026}
		\engtaskenddate{17}{6}{2026}
		\engtaskdescription{Implementar el repositorio de restricciones con filtros por versión y estado, búsqueda paginada y recuperación de expresiones normalizadas para procesos de validación lógica.}
		\engtaskprogrammer{Jose Luis Echemendia López}
		\end{engineeringtask}

		\begin{engineeringtask}[tb:engtask_t044]
		\engtaskuserstory{9}
		\engtaskname{Implementar servicio de restricciones transversales}
		\engtasktype{Crear servicio}
		\engtaskpointestimation{12}
		\engtaskstartdate{17}{6}{2026}
		\engtaskenddate{18}{6}{2026}
		\engtaskdescription{Implementar el servicio de restricciones transversales para validar sintaxis, transformar expresiones a forma operativa y preservar trazabilidad de cambios mediante reglas de versionado.}
		\engtaskprogrammer{Jose Luis Echemendia López}
		\end{engineeringtask}

		\begin{engineeringtask}[tb:engtask_t045]
		\engtaskuserstory{9}
		\engtaskname{Exponer endpoints de restricciones}
		\engtasktype{Crear controlador}
		\engtaskpointestimation{8}
		\engtaskstartdate{18}{6}{2026}
		\engtaskenddate{19}{6}{2026}
		\engtaskdescription{Implementar el controlador de restricciones lógicas para alta, consulta, actualización y cambio de estado, asegurando validación sintáctica de expresiones, consistencia de versión y respuestas de error homogéneas para fallos de lógica.}
		\engtaskprogrammer{Jose Luis Echemendia López}
		\end{engineeringtask}

		\subsubsection*{HU-10: Integración y Sincronización UVL}

		\begin{engineeringtask}[tb:engtask_t046]
		\engtaskuserstory{10}
		\engtaskname{Implementar importador UVL a estructura}
		\engtasktype{Crear clase}
		\engtaskpointestimation{14}
		\engtaskstartdate{22}{6}{2026}
		\engtaskenddate{23}{6}{2026}
		\engtaskdescription{Diseñar e implementar la clase FeatureModelUVLImporter para parseo, validación semántica y construcción de estructura interna a partir de UVL, con reporte de errores orientado a diagnóstico técnico.}
		\engtaskprogrammer{Jose Luis Echemendia López}
		\end{engineeringtask}

		\begin{engineeringtask}[tb:engtask_t047]
		\engtaskuserstory{10}
		\engtaskname{Implementar esquemas de sincronización UVL}
		\engtasktype{Crear esquemas de validación entrada de la api y salida}
		\engtaskpointestimation{6}
		\engtaskstartdate{23}{6}{2026}
		\engtaskenddate{23}{6}{2026}
		\engtaskdescription{Definir esquemas de entrada y salida para persistencia de UVL, sincronización desde estructura y aplicación sobre versión destino, con validaciones de compatibilidad y tamaño de contenido.}
		\engtaskprogrammer{Jose Luis Echemendia López}
		\end{engineeringtask}

		\begin{engineeringtask}[tb:engtask_t048]
		\engtaskuserstory{10}
		\engtaskname{Implementar servicio de sincronización UVL}
		\engtasktype{Crear servicio}
		\engtaskpointestimation{12}
		\engtaskstartdate{23}{6}{2026}
		\engtaskenddate{24}{6}{2026}
		\engtaskdescription{Implementar el servicio de sincronización UVL para garantizar equivalencia entre representación visual y textual, persistir UVL efectivo y gestionar creación controlada de nuevas versiones.}
		\engtaskprogrammer{Jose Luis Echemendia López}
		\end{engineeringtask}

		\begin{engineeringtask}[tb:engtask_t049]
		\engtaskuserstory{10}
		\engtaskname{Exponer endpoints de integración UVL}
		\engtasktype{Crear controlador}
		\engtaskpointestimation{8}
		\engtaskstartdate{24}{6}{2026}
		\engtaskenddate{25}{6}{2026}
		\engtaskdescription{Implementar el controlador de integración UVL para guardar contenido validado, sincronizar UVL desde estructura y aplicar UVL sobre una versión destino; contemplar validaciones previas de consistencia y mensajes de error orientados a diagnóstico.}
		\engtaskprogrammer{Jose Luis Echemendia López}
		\end{engineeringtask}

		\begin{engineeringtask}[tb:engtask_t050]
		\engtaskuserstory{10}
		\engtaskname{Definir excepciones de UVL}
		\engtasktype{Crear excepciones personalizadas}
		\engtaskpointestimation{4}
		\engtaskstartdate{25}{6}{2026}
		\engtaskenddate{26}{6}{2026}
		\engtaskdescription{Definir excepciones especializadas para errores de parseo, inconsistencias de sincronización y conflictos estructurales, estandarizando el reporte técnico para depuración y soporte.}
		\engtaskprogrammer{Jose Luis Echemendia López}
		\end{engineeringtask}

		\subsubsection*{HU-11: Validación Lógica y Estructural del Modelo}

		\begin{engineeringtask}[tb:engtask_t051]
		\engtaskuserstory{11}
		\engtaskname{Implementar validador lógico SAT/SMT}
		\engtasktype{Crear clase}
		\engtaskpointestimation{16}
		\engtaskstartdate{29}{6}{2026}
		\engtaskenddate{30}{6}{2026}
		\engtaskdescription{Diseñar e implementar la clase FeatureModelLogicalValidator para verificar satisfacibilidad global, consistencia lógica y validez de configuraciones propuestas sobre un modelo y versión concretos.}
		\engtaskprogrammer{Jose Luis Echemendia López}
		\end{engineeringtask}

		\begin{engineeringtask}[tb:engtask_t052]
		\engtaskuserstory{11}
		\engtaskname{Implementar analizador estructural}
		\engtasktype{Crear clase}
		\engtaskpointestimation{12}
		\engtaskstartdate{30}{6}{2026}
		\engtaskenddate{1}{7}{2026}
		\engtaskdescription{Diseñar e implementar el analizador estructural para detección de ciclos, nodos huérfanos, características muertas, redundancias y componentes fuertemente conectados.}
		\engtaskprogrammer{Jose Luis Echemendia López}
		\end{engineeringtask}

		\begin{engineeringtask}[tb:engtask_t053]
		\engtaskuserstory{11}
		\engtaskname{Implementar fachada de análisis del FM}
		\engtasktype{Crear servicio}
		\engtaskpointestimation{10}
		\engtaskstartdate{1}{7}{2026}
		\engtaskenddate{2}{7}{2026}
		\engtaskdescription{Implementar la fachada de análisis para orquestar validaciones lógicas y estructurales, consolidar resultados en un único contrato y normalizar severidad, trazas y recomendaciones técnicas.}
		\engtaskprogrammer{Jose Luis Echemendia López}
		\end{engineeringtask}

		\begin{engineeringtask}[tb:engtask_t054]
		\engtaskuserstory{11}
		\engtaskname{Implementar esquemas de respuesta de análisis}
		\engtasktype{Crear esquemas de validación entrada de la api y salida}
		\engtaskpointestimation{6}
		\engtaskstartdate{2}{7}{2026}
		\engtaskenddate{2}{7}{2026}
		\engtaskdescription{Definir esquemas de respuesta para reportes de análisis lógico y estructural, incluyendo severidad, evidencia, recomendaciones y metadatos de ejecución para consumo por clientes y observabilidad.}
		\engtaskprogrammer{Jose Luis Echemendia López}
		\end{engineeringtask}

		\begin{engineeringtask}[tb:engtask_t055]
		\engtaskuserstory{11}
		\engtaskname{Exponer endpoints de validación y análisis}
		\engtasktype{Crear controlador}
		\engtaskpointestimation{8}
		\engtaskstartdate{2}{7}{2026}
		\engtaskenddate{3}{7}{2026}
		\engtaskdescription{Implementar el controlador de validación para análisis lógico y estructural del modelo, validación de configuraciones propuestas y reporte de hallazgos; incorporar ejecución segura con límites de tiempo y manejo controlado de excepciones de análisis.}
		\engtaskprogrammer{Jose Luis Echemendia López}
		\end{engineeringtask}

		\subsubsection*{HU-12: Gestión de Configuraciones}

		\begin{engineeringtask}[tb:engtask_t056]
		\engtaskuserstory{12}
		\engtaskname{Ajustar modelo de Configuration}
		\engtasktype{Definir modelos}
		\engtaskpointestimation{8}
		\engtaskstartdate{6}{7}{2026}
		\engtaskenddate{6}{7}{2026}
		\engtaskdescription{Definir y ajustar el modelo de Configuration con vínculo explícito a versión objetivo, selección de características, estado operativo e historial de cambios para preservar validez temporal.}
		\engtaskprogrammer{Jose Luis Echemendia López}
		\end{engineeringtask}

		\begin{engineeringtask}[tb:engtask_t057]
		\engtaskuserstory{12}
		\engtaskname{Implementar esquemas de configuraciones persistentes}
		\engtasktype{Crear esquemas de validación entrada de la api y salida}
		\engtaskpointestimation{8}
		\engtaskstartdate{6}{7}{2026}
		\engtaskenddate{7}{7}{2026}
		\engtaskdescription{Definir esquemas para creación, consulta, listado, actualización y cambio de estado de configuraciones persistentes, así como validación de propuestas no persistidas con contratos de salida normalizados.}
		\engtaskprogrammer{Jose Luis Echemendia López}
		\end{engineeringtask}

		\begin{engineeringtask}[tb:engtask_t058]
		\engtaskuserstory{12}
		\engtaskname{Implementar repositorio de configuraciones}
		\engtasktype{Crear repositorio}
		\engtaskpointestimation{8}
		\engtaskstartdate{7}{7}{2026}
		\engtaskenddate{8}{7}{2026}
		\engtaskdescription{Implementar el repositorio de configuraciones con persistencia paginada, filtros por versión y estado, y carga eficiente de asociaciones con características seleccionadas.}
		\engtaskprogrammer{Jose Luis Echemendia López}
		\end{engineeringtask}

		\begin{engineeringtask}[tb:engtask_t059]
		\engtaskuserstory{12}
		\engtaskname{Implementar servicio de gestión de configuraciones}
		\engtasktype{Crear servicio}
		\engtaskpointestimation{12}
		\engtaskstartdate{8}{7}{2026}
		\engtaskenddate{9}{7}{2026}
		\engtaskdescription{Implementar el servicio de configuraciones para validar consistencia funcional, persistir selecciones por versión objetivo y gestionar reglas de activación/desactivación con control de estado.}
		\engtaskprogrammer{Jose Luis Echemendia López}
		\end{engineeringtask}

		\begin{engineeringtask}[tb:engtask_t060]
		\engtaskuserstory{12}
		\engtaskname{Exponer endpoints de configuraciones}
		\engtasktype{Crear controlador}
		\engtaskpointestimation{8}
		\engtaskstartdate{9}{7}{2026}
		\engtaskenddate{10}{7}{2026}
		\engtaskdescription{Implementar el controlador de configuraciones para creación persistente, consulta por identificador, listado paginado, actualización y activación/desactivación; incluir endpoint de validación de selección no persistente con contratos claros de entrada y salida.}
		\engtaskprogrammer{Jose Luis Echemendia López}
		\end{engineeringtask}

		\subsubsection*{HU-13: Validación y Generación Automática de Configuraciones}

		\begin{engineeringtask}[tb:engtask_t061]
		\engtaskuserstory{13}
		\engtaskname{Implementar núcleo de generación de configuraciones}
		\engtasktype{Crear clase}
		\engtaskpointestimation{10}
		\engtaskstartdate{13}{7}{2026}
		\engtaskenddate{14}{7}{2026}
		\engtaskdescription{Diseñar e implementar la clase base de generación de configuraciones válidas, incorporando estrategias iniciales deterministas y aleatorias, validación previa de restricciones y criterios de factibilidad sobre la versión objetivo del modelo.}
		\engtaskprogrammer{Jose Luis Echemendia López}
		\end{engineeringtask}

		\begin{engineeringtask}[tb:engtask_t062]
		\engtaskuserstory{13}
		\engtaskname{Implementar estrategia BEAM para generación válida}
		\engtasktype{Crear clase}
		\engtaskpointestimation{8}
		\engtaskstartdate{14}{7}{2026}
		\engtaskenddate{14}{7}{2026}
		\engtaskdescription{Extender el generador con búsqueda BEAM, incluyendo funciones de puntuación, poda y control de diversidad, para producir candidatos válidos con mejor cobertura y calidad estructural.}
		\engtaskprogrammer{Jose Luis Echemendia López}
		\end{engineeringtask}

		\begin{engineeringtask}[tb:engtask_t063]
		\engtaskuserstory{13}
		\engtaskname{Implementar optimización mono-objetivo}
		\engtasktype{Crear servicio}
		\engtaskpointestimation{8}
		\engtaskstartdate{14}{7}{2026}
		\engtaskenddate{15}{7}{2026}
		\engtaskdescription{Implementar el servicio de optimización mono-objetivo con algoritmos genéticos, definiendo función de aptitud, operadores evolutivos y criterios de convergencia para seleccionar configuraciones válidas de mayor calidad.}
		\engtaskprogrammer{Jose Luis Echemendia López}
		\end{engineeringtask}

		\begin{engineeringtask}[tb:engtask_t064]
		\engtaskuserstory{13}
		\engtaskname{Implementar optimización multiobjetivo}
		\engtasktype{Crear servicio}
		\engtaskpointestimation{8}
		\engtaskstartdate{15}{7}{2026}
		\engtaskenddate{15}{7}{2026}
		\engtaskdescription{Implementar el servicio de optimización multiobjetivo con enfoque de frentes de Pareto, equilibrando criterios de calidad contrapuestos y garantizando diversidad y estabilidad en el conjunto de soluciones candidatas.}
		\engtaskprogrammer{Jose Luis Echemendia López}
		\end{engineeringtask}

		\begin{engineeringtask}[tb:engtask_t065]
		\engtaskuserstory{13}
		\engtaskname{Implementar esquemas de generación/optimización}
		\engtasktype{Crear esquemas de validación entrada de la api y salida}
		\engtaskpointestimation{6}
		\engtaskstartdate{15}{7}{2026}
		\engtaskenddate{16}{7}{2026}
		\engtaskdescription{Definir contratos de API para validación de propuestas, generación automática y optimización, incluyendo estrategia seleccionada, parámetros de ejecución, límites operativos y estructura homogénea de resultados.}
		\engtaskprogrammer{Jose Luis Echemendia López}
		\end{engineeringtask}

		\begin{engineeringtask}[tb:engtask_t066]
		\engtaskuserstory{13}
		\engtaskname{Exponer endpoints de validación y generación automática}
		\engtasktype{Crear controlador}
		\engtaskpointestimation{8}
		\engtaskstartdate{16}{7}{2026}
		\engtaskenddate{17}{7}{2026}
		\engtaskdescription{Implementar el controlador de configuración avanzada para validar propuestas, generar configuraciones y ejecutar optimización con parametrización controlada; estandarizar respuestas y trazas de ejecución para consumo cliente.}
		\engtaskprogrammer{Jose Luis Echemendia López}
		\end{engineeringtask}

		\begin{engineeringtask}[tb:engtask_t067]
		\engtaskuserstory{13}
		\engtaskname{Ejecutar optimización intensiva en background}
		\engtasktype{Crear tareas de celery}
		\engtaskpointestimation{8}
		\engtaskstartdate{17}{7}{2026}
		\engtaskenddate{18}{7}{2026}
		\engtaskdescription{Implementar flujo asíncrono de optimización intensiva con colas de trabajo, control de reintentos, persistencia de estado, almacenamiento de resultados y reanudación por identificador de ejecución.}
		\engtaskprogrammer{Jose Luis Echemendia López}
		\end{engineeringtask}

		\subsubsection*{HU-14: Enriquecimiento con Recursos y Taxonomías}

		\begin{engineeringtask}[tb:engtask_t068]
		\engtaskuserstory{14}
		\engtaskname{Ajustar modelos de Tag y Resource}
		\engtasktype{Definir modelos}
		\engtaskpointestimation{10}
		\engtaskstartdate{20}{7}{2026}
		\engtaskenddate{21}{7}{2026}
		\engtaskdescription{Definir y ajustar los modelos de Tag y Resource con reglas de unicidad, metadatos extendidos y asociaciones con características y propietario, garantizando consistencia funcional y persistencia estable.}
		\engtaskprogrammer{Jose Luis Echemendia López}
		\end{engineeringtask}

		\begin{engineeringtask}[tb:engtask_t069]
		\engtaskuserstory{14}
		\engtaskname{Implementar esquemas de etiquetas y recursos}
		\engtasktype{Crear esquemas de validación entrada de la api y salida}
		\engtaskpointestimation{8}
		\engtaskstartdate{21}{7}{2026}
		\engtaskenddate{21}{7}{2026}
		\engtaskdescription{Definir esquemas para gestión completa de etiquetas y gestión parcial/completa de recursos, además de operaciones de asociación y desasociación con características, con validaciones de negocio y formato.}
		\engtaskprogrammer{Jose Luis Echemendia López}
		\end{engineeringtask}

		\begin{engineeringtask}[tb:engtask_t070]
		\engtaskuserstory{14}
		\engtaskname{Implementar repositorios de taxonomías y recursos}
		\engtasktype{Crear repositorio}
		\engtaskpointestimation{10}
		\engtaskstartdate{21}{7}{2026}
		\engtaskenddate{22}{7}{2026}
		\engtaskdescription{Implementar repositorios de taxonomías y recursos con búsquedas paginadas, filtros por metadatos y operaciones de vinculación entre entidades, asegurando integridad relacional y consistencia de lectura.}
		\engtaskprogrammer{Jose Luis Echemendia López}
		\end{engineeringtask}

		\begin{engineeringtask}[tb:engtask_t071]
		\engtaskuserstory{14}
		\engtaskname{Implementar servicio de enriquecimiento semántico}
		\engtasktype{Crear servicio}
		\engtaskpointestimation{12}
		\engtaskstartdate{22}{7}{2026}
		\engtaskenddate{23}{7}{2026}
		\engtaskdescription{Implementar el servicio de enriquecimiento semántico para aplicar reglas de propiedad, versionado de metadatos y consistencia de asociaciones entre etiquetas, características y recursos educativos.}
		\engtaskprogrammer{Jose Luis Echemendia López}
		\end{engineeringtask}

		\begin{engineeringtask}[tb:engtask_t072]
		\engtaskuserstory{14}
		\engtaskname{Exponer endpoints de tags y recursos}
		\engtasktype{Crear controlador}
		\engtaskpointestimation{8}
		\engtaskstartdate{23}{7}{2026}
		\engtaskenddate{24}{7}{2026}
		\engtaskdescription{Implementar el controlador de taxonomías y recursos para gestión de etiquetas, asociación/desasociación con características y administración de recursos educativos con metadatos; aplicar reglas de propietario, validaciones de negocio y respuestas consistentes para operaciones parciales y completas.}
		\engtaskprogrammer{Jose Luis Echemendia López}
		\end{engineeringtask}
		
	\chapter{Pruebas de aceptación}
	\label{anx:pa}

		\subsubsection*{Épica 1: Gestión de Identidad y Acceso}
		\subsubsection*{HU-01: Autenticación y Autogestión de cuenta}
		
		\begin{acceptancetest}[tb:HU1_P1]
		\testcasecode{HU1\_P1}
		\testcasename{Registro de usuario con correo y contraseña}
		\testcasedescription{Validar que un usuario nuevo pueda registrarse en el sistema con correo y contraseña válidos.}
		\testcaseexeccond{
		- El servidor backend está disponible.\\
		- La base de datos está accesible.\\
		- No existe un usuario registrado con el correo proporcionado.
		}
		\testcaseexecstep{
		1. Enviar POST a /api/v1/auth/register con payload: \{"email": "usuario@test.com", "password": "Pass123!@"\}.\\
		2. Validar que la respuesta contenga status 201.\\
		3. Validar que la respuesta incluya el token de autenticación.\\
		4. Validar que se creó el usuario en la base de datos.
		}
		\testcaseexpresult{
		- Status HTTP: 201 Created.\\
		- Respuesta contiene access\_token, refresh\_token y user\_id.\\
		- Usuario registrado en BD con is\_active=True.\\
		- Contraseña almacenada como hash (no en texto plano).
		}
		\testcaseuserstory{1}
		\end{acceptancetest}

		\begin{acceptancetest}[tb:HU1_P2]
		\testcasecode{HU1\_P2}
		\testcasename{Inicio de sesión con correo y contraseña}
		\testcasedescription{Validar que un usuario registrado pueda iniciar sesión correctamente.}
		\testcaseexeccond{
		- El servidor backend está disponible.\\
		- Existe un usuario registrado con correo test@example.com y contraseña conocida.
		}
		\testcaseexecstep{
		1. Enviar POST a /api/v1/auth/login con payload: \{"email": "test@example.com", "password": "Pass123!@"\}.\\
		2. Validar que la respuesta contenga status 200.\\
		3. Verificar que se devuelve un access\_token válido.
		}
		\testcaseexpresult{
		- Status HTTP: 200 OK.\\
		- Respuesta contiene access\_token, refresh\_token y user\_id.\\
		- El token es válido para futuras peticiones autenticadas.
		}
		\testcaseuserstory{1}
		\end{acceptancetest}

		\begin{acceptancetest}[tb:HU1_P3]
		\testcasecode{HU1\_P3}
		\testcasename{Recuperación de acceso mediante token por correo}
		\testcasedescription{Validar que un usuario pueda solicitar recuperación de contraseña y recibir token por correo.}
		\testcaseexeccond{
		- El servidor backend está disponible.\\
		- El servicio de correo está configurado.\\
		- Existe un usuario con correo usuario@test.com.\\
		- El usuario se ha registrado previamente.
		}
		\testcaseexecstep{
		1. Enviar POST a /api/v1/auth/forgot-password con payload: \{"email": "usuario@test.com"\}.\\
		2. Validar que la respuesta contenga status 200.\\
		3. Verificar que se envió un correo con token de recuperación.\\
		4. Extraer token del correo y enviar PUT a /api/v1/auth/reset-password con payload: \{"token": "<token>", "new\_password": "NewPass456!@"\}.\\
		5. Validar que la contraseña se actualizó.
		}
		\testcaseexpresult{
		- Status HTTP: 200 OK (solicitud).\\
		- Se envía correo con token de recuperación.\\
		- Status HTTP: 200 OK (reset).\\
		- Usuario puede iniciar sesión con nueva contraseña.\\
		- Token anterior se invalida.
		}
		\testcaseuserstory{1}
		\end{acceptancetest}

		\begin{acceptancetest}[tb:HU1_P4]
		\testcasecode{HU1\_P4}
		\testcasename{Edición de perfil propio sin intervención administrativa}
		\testcasedescription{Validar que un usuario autenticado pueda actualizar su perfil (correo, nombre) sin necesidad de administrador.}
		\testcaseexeccond{
		- Usuario está autenticado con token válido.\\
		- El nuevo correo no está registrado en otro usuario.\\
		- El usuario tiene permisos para editar su propio perfil.
		}
		\testcaseexecstep{
		1. Enviar PUT a /api/v1/auth/profile con headers de autenticación y payload: \{"email": "nuevo@test.com", "first\_name": "Juan", "last\_name": "Pérez"\}.\\
		2. Validar que la respuesta contenga status 200.\\
		3. Consultar GET /api/v1/auth/profile para confirmar cambios.\\
		4. Validar que los datos actualizados coinciden.
		}
		\testcaseexpresult{
		- Status HTTP: 200 OK.\\
		- Perfil actualizado en base de datos.\\
		- Respuesta contiene los datos actualizados.\\
		- Token sigue siendo válido.\\
		- Cambios persistidos en BD.
		}
		\testcaseuserstory{1}
		\end{acceptancetest}

		\begin{acceptancetest}[tb:HU1_P5]
		\testcasecode{HU1\_P5}
		\testcasename{Cierre de sesión}
		\testcasedescription{Validar que un usuario pueda cerrar sesión invalidando su token.}
		\testcaseexeccond{
		- Usuario está autenticado con token válido.\\
		- El servidor mantiene una lista de tokens revocados.
		}
		\testcaseexecstep{
		1. Enviar POST a /api/v1/auth/logout con headers de autenticación.\\
		2. Validar que la respuesta contenga status 200.\\
		3. Intentar acceder a un endpoint protegido con el token anterior.\\
		4. Validar que se rechaza la petición.
		}
		\testcaseexpresult{
		- Status HTTP: 200 OK.\\
		- Token se añade a lista de revocación.\\
		- Peticiones posteriores con ese token retornan 401 Unauthorized.\\
		- Mensaje: "Token revoked or expired".
		}
		\testcaseuserstory{1}
		\end{acceptancetest}

		\subsubsection*{HU-02: Administración de Usuarios y Permisos}
	
		\begin{acceptancetest}[tb:HU2_P1]
		\testcasecode{HU2\_P1}
		\testcasename{Listar usuarios con filtro por rol}
		\testcasedescription{Validar que un administrador puede obtener listado de usuarios filtrados por rol.}
		\testcaseexeccond{
		- Usuario es administrador (rol=ADMIN).\\
		- Existen usuarios con diferentes roles en la BD.\\
		- Usuario está autenticado con token de administrador.
		}
		\testcaseexecstep{
		1. Enviar GET a /api/v1/admin/users?role=EDITOR\&skip=0\&limit=10 con headers de autenticación.\\
		2. Validar que la respuesta contenga status 200.\\
		3. Verificar que todos los usuarios retornados tienen role="EDITOR".\\
		4. Validar que la respuesta incluye total, skip, limit.
		}
		\testcaseexpresult{
		- Status HTTP: 200 OK.\\
		- Respuesta contiene arreglo de usuarios filtrados por rol.\\
		- Cada usuario incluye: id, email, first\_name, last\_name, role, is\_active.\\
		- Paginación correcta.
		}
		\testcaseuserstory{2}
		\end{acceptancetest}

		\begin{acceptancetest}[tb:HU2_P2]
		\testcasecode{HU2\_P2}
		\testcasename{Listar usuarios con filtro por nombre}
		\testcasedescription{Validar que un administrador puede filtrar usuarios por nombre completo.}
		\testcaseexeccond{
		- Usuario es administrador.\\
		- Existen usuarios registrados con nombres distintos.\\
		- Usuario está autenticado.
		}
		\testcaseexecstep{
		1. Enviar GET a /api/v1/admin/users?name=Juan\&skip=0\&limit=10 con headers de autenticación.\\
		2. Validar que la respuesta contenga status 200.\\
		3. Verificar que todos los usuarios retornados contienen "Juan" en nombre o apellido.
		}
		\testcaseexpresult{
		- Status HTTP: 200 OK.\\
		- Respuesta contiene usuarios cuyo nombre/apellido contiene la búsqueda.\\
		- Búsqueda es case-insensitive.
		}
		\testcaseuserstory{2}
		\end{acceptancetest}

		\begin{acceptancetest}[tb:HU2_P3]
		\testcasecode{HU2\_P3}
		\testcasename{Modificar rol de usuario}
		\testcasedescription{Validar que un administrador puede cambiar el rol asignado a un usuario.}
		\testcaseexeccond{
		- Usuario es administrador.\\
		- Existe un usuario con rol VIEWER.\\
		- Usuario está autenticado con credenciales de admin.
		}
		\testcaseexecstep{
		1. Obtener ID de usuario a modificar.\\
		2. Enviar PUT a /api/v1/admin/users/\{user\_id\}/role con payload: \{"role": "EDITOR"\}.\\
		3. Validar que la respuesta contenga status 200.\\
		4. Consultar GET /api/v1/admin/users/\{user\_id\} para confirmar cambio.\\
		5. Verificar que role="EDITOR".
		}
		\testcaseexpresult{
		- Status HTTP: 200 OK.\\
		- Respuesta contiene usuario actualizado con nuevo rol.\\
		- BD refleja el cambio.\\
		- Permisos del usuario se actualizan dinámicamente.
		}
		\testcaseuserstory{2}
		\end{acceptancetest}

		\begin{acceptancetest}[tb:HU2_P4]
		\testcasecode{HU2\_P4}
		\testcasename{Desactivar cuenta de usuario sin eliminar información}
		\testcasedescription{Validar que un administrador puede desactivar una cuenta sin eliminar los datos.}
		\testcaseexeccond{
		- Usuario es administrador.\\
		- Existe un usuario activo.\\
		- No se permiten borrados físicos en el sistema.
		}
		\testcaseexecstep{
		1. Obtener ID de usuario a desactivar.\\
		2. Enviar PUT a /api/v1/admin/users/\{user\_id\}/status con payload: \{"is\_active": false\}.\\
		3. Validar que la respuesta contenga status 200.\\
		4. Intentar que el usuario desactivado inicie sesión.\\
		5. Verificar que el login falla con mensaje de cuenta desactivada.
		}
		\testcaseexpresult{
		- Status HTTP: 200 OK.\\
		- Usuario marcado como is\_active=False en BD.\\
		- Datos del usuario persisten en BD.\\
		- Usuario no puede autenticarse.\\
		- Mensaje: "User account is deactivated".
		}
		\testcaseuserstory{2}
		\end{acceptancetest}

		\begin{acceptancetest}[tb:HU2_P5]
		\testcasecode{HU2\_P5}
		\testcasename{Reactivar cuenta de usuario}
		\testcasedescription{Validar que un administrador puede reactivar una cuenta desactivada.}
		\testcaseexeccond{
		- Usuario es administrador.\\
		- Existe un usuario desactivado (is\_active=False).\\
		- Los datos del usuario están intactos en BD.
		}
		\testcaseexecstep{
		1. Obtener ID de usuario a reactivar.\\
		2. Enviar PUT a /api/v1/admin/users/\{user\_id\}/status con payload: \{"is\_active": true\}.\\
		3. Validar que la respuesta contenga status 200.\\
		4. Intentar autenticar con el usuario reactivado.\\
		5. Verificar que el login es exitoso.
		}
		\testcaseexpresult{
		- Status HTTP: 200 OK.\\
		- Usuario marcado como is\_active=True.\\
		- Datos originales restaurados.\\
		- Usuario puede autenticarse normalmente.
		}
		\testcaseuserstory{2}
		\end{acceptancetest}

		\subsubsection*{Épica 2: Gestión de Dominios Curriculares}
		\subsubsection*{HU-03: Gestión de Dominios de Conocimiento}

		\begin{acceptancetest}[tb:HU3_P1]
		\testcasecode{HU3\_P1}
		\testcasename{Crear dominio con nombre y descripción}
		\testcasedescription{Validar que un administrador puede crear un nuevo dominio curricular.}
		\testcaseexeccond{
		- Usuario es administrador.\\
		- Usuario está autenticado.\\
		- No existe un dominio con el mismo nombre.
		}
		\testcaseexecstep{
		1. Enviar POST a /api/v1/domains con payload: \{"name": "Ingeniería de Software", "description": "Dominio para estudios en ingeniería de software"\}.\\
		2. Validar que la respuesta contenga status 201.\\
		3. Verificar que se retorna id, name, description, is\_active, created\_at.\\
		4. Consultar BD para confirmar persistencia.
		}
		\testcaseexpresult{
		- Status HTTP: 201 Created.\\
		- Respuesta contiene ID generado y campos completos.\\
		- Dominio se crea con is\_active=True por defecto.\\
		- created\_at y updated\_at se establecen automáticamente.
		}
		\testcaseuserstory{3}
		\end{acceptancetest}

		\begin{acceptancetest}[tb:HU3_P2]
		\testcasecode{HU3\_P2}
		\testcasename{Obtener dominio con modelos asociados}
		\testcasedescription{Validar que se puede obtener un dominio junto con todos sus modelos vinculados.}
		\testcaseexeccond{
		- Existe un dominio con ID conocido.\\
		- El dominio tiene modelos asociados.\\
		- Usuario está autenticado.
		}
		\testcaseexecstep{
		1. Enviar GET a /api/v1/domains/\{domain\_id\}?include\_models=true con headers de autenticación.\\
		2. Validar que la respuesta contenga status 200.\\
		3. Verificar que la respuesta incluye id, name, description, is\_active.\\
		4. Verificar que incluye array models con todos los modelos del dominio.\\
		5. Cada modelo debe contener: id, name, domain\_id, created\_at.
		}
		\testcaseexpresult{
		- Status HTTP: 200 OK.\\
		- Respuesta contiene información del dominio.\\
		- Array models contiene todos los modelos vinculados.\\
		- Modelos activos e inactivos se incluyen (según filtro).
		}
		\testcaseuserstory{3}
		\end{acceptancetest}

		\begin{acceptancetest}[tb:HU3_P3]
		\testcasecode{HU3\_P3}
		\testcasename{Desactivar dominio}
		\testcasedescription{Validar que un administrador puede desactivar un dominio sin eliminar información.}
		\testcaseexeccond{
		- Usuario es administrador.\\
		- Existe un dominio activo.\\
		- El dominio puede no tener modelos o tenerlos.
		}
		\testcaseexecstep{
		1. Obtener ID de dominio a desactivar.\\
		2. Enviar PUT a /api/v1/domains/\{domain\_id\} con payload: \{"is\_active": false\}.\\
		3. Validar que la respuesta contenga status 200.\\
		4. Consultar GET /api/v1/domains/\{domain\_id\} para confirmar.\\
		5. Verificar que is\_active=False.
		}
		\testcaseexpresult{
		- Status HTTP: 200 OK.\\
		- Dominio marcado como is\_active=False.\\
		- Información persiste en BD.\\
		- Modelos del dominio no son eliminados.
		}
		\testcaseuserstory{3}
		\end{acceptancetest}

		\begin{acceptancetest}[tb:HU3_P4]
		\testcasecode{HU3\_P4}
		\testcasename{Listar todos los dominios}
		\testcasedescription{Validar que se puede obtener un listado paginado de dominios.}
		\testcaseexeccond{
		- Existen múltiples dominios en la BD.\\
		- Usuario está autenticado.
		}
		\testcaseexecstep{
		1. Enviar GET a /api/v1/domains?skip=0\&limit=10 con headers de autenticación.\\
		2. Validar que la respuesta contenga status 200.\\
		3. Verificar que se retorna array de dominios.\\
		4. Cada dominio incluye: id, name, description, is\_active, created\_at.\\
		5. Validar paginación: total, skip, limit.
		}
		\testcaseexpresult{
		- Status HTTP: 200 OK.\\
		- Array de dominios ordenados por fecha.\\
		- Paginación correcta.\\
		- Total de registros se incluye.
		}
		\testcaseuserstory{3}
		\end{acceptancetest}

		\subsubsection*{Épica 3: Ciclo de Vida del \ac{FM}}
		\subsubsection*{HU-04: Gestión Operativa de \ac{FM}}
		
		\begin{acceptancetest}[tb:HU4_P1]
		\testcasecode{HU4\_P1}
		\testcasename{Crear modelo asociado a dominio válido}
		\testcasedescription{Validar que se puede crear un nuevo modelo de características en un dominio existente.}
		\testcaseexeccond{
		- Usuario es diseñador o administrador.\\
		- Existe un dominio activo.\\
		- Usuario está autenticado.\\
		- No existe modelo con igual nombre en el dominio.
		}
		\testcaseexecstep{
		1. Obtener ID de dominio válido.\\
		2. Enviar POST a /api/v1/models con payload: \{"name": "Plan 2024", "description": "Modelo curricular 2024", "domain\_id": "<domain\_id>"\}.\\
		3. Validar que la respuesta contenga status 201.\\
		4. Verificar que se retorna id, name, domain\_id, created\_at.\\
		5. Validar que existe en BD.
		}
		\testcaseexpresult{
		- Status HTTP: 201 Created.\\
		- Modelo creado con is\_active=True.\\
		- Respuesta incluye ID generado.\\
		- domain\_id está vinculado correctamente.\\
		- Timestamps (created\_at, updated\_at) se establecen.
		}
		\testcaseuserstory{4}
		\end{acceptancetest}

		\begin{acceptancetest}[tb:HU4_P2]
		\testcasecode{HU4\_P2}
		\testcasename{Actualizar metadatos del modelo}
		\testcasedescription{Validar que los metadatos de un modelo pueden actualizarse sin afectar la estructura.}
		\testcaseexeccond{
		- Existe un modelo con ID conocido.\\
		- Usuario es diseñador del modelo o administrador.\\
		- Usuario está autenticado.
		}
		\testcaseexecstep{
		1. Obtener ID de modelo.\\
		2. Enviar PUT a /api/v1/models/\{model\_id\} con payload: \{"name": "Plan 2024 v2", "description": "Modelo actualizado"\}.\\
		3. Validar que la respuesta contenga status 200.\\
		4. Consultar GET /api/v1/models/\{model\_id\} para confirmar.\\
		5. Verificar que metadata está actualizada y estructura intacta.
		}
		\testcaseexpresult{
		- Status HTTP: 200 OK.\\
		- Metadatos actualizados en BD.\\
		- Estructura del modelo no se modifica.\\
		- updated\_at se actualiza.\\
		- Versiones no se ven afectadas.
		}
		\testcaseuserstory{4}
		\end{acceptancetest}

		\begin{acceptancetest}[tb:HU4_P3]
		\testcasecode{HU4\_P3}
		\testcasename{Listar modelos con paginación y filtrado}
		\testcasedescription{Validar que se pueden listar modelos con paginación, ordenamiento y filtros.}
		\testcaseexeccond{
		- Existen múltiples modelos en la BD.\\
		- Modelos están vinculados a dominios.\\
		- Usuario está autenticado.
		}
		\testcaseexecstep{
		1. Enviar GET a /api/v1/models?domain\_id=<id>\&skip=0\&limit=10\&sort=-created\_at con headers.\\
		2. Validar que la respuesta contenga status 200.\\
		3. Verificar que se retorna array de modelos.\\
		4. Validar paginación: total, skip, limit.\\
		5. Verificar que ordenamiento es correcto (descendente por created\_at).
		}
		\testcaseexpresult{
		- Status HTTP: 200 OK.\\
		- Array de modelos paginados.\\
		- Filtrado por dominio correctamente aplicado.\\
		- Ordenamiento correcto.\\
		- Total de registros incluido.
		}
		\testcaseuserstory{4}
		\end{acceptancetest}

		\subsubsection*{HU-05: Gestión de Versiones y Publicación}

		\begin{acceptancetest}[tb:HU5_P1]
		\testcasecode{HU5\_P1}
		\testcasename{Crear nueva versión mediante copia al escribir}
		\testcasedescription{Validar que se puede crear una nueva versión duplicando una versión existente (Copy-on-Write).}
		\testcaseexeccond{
		- Existe un modelo con versión publicada.\\
		- Usuario es diseñador del modelo.\\
		- Usuario está autenticado.
		}
		\testcaseexecstep{
		1. Obtener ID de modelo e ID de versión origen.\\
		2. Enviar POST a /api/v1/models/\{model\_id\}/versions con payload: \{"version\_origin\_id": "<version\_id>", "description": "Version 2.0"\}.\\
		3. Validar que la respuesta contenga status 201.\\
		4. Verificar que se crea versión con estado "DRAFT".\\
		5. Validar que la estructura se copió desde versión origen.
		}
		\testcaseexpresult{
		- Status HTTP: 201 Created.\\
		- Nueva versión creada con estado DRAFT.\\
		- Copia completa de la estructura anterior.\\
		- Versión origen no se modifica.\\
		- Versión nueva tiene version\_number incrementado.
		}
		\testcaseuserstory{5}
		\end{acceptancetest}

		\begin{acceptancetest}[tb:HU5_P2]
		\testcasecode{HU5\_P2}
		\testcasename{Validar versión antes de publicación}
		\testcasedescription{Validar que se ejecutan validaciones antes de permitir publicación de versión.}
		\testcaseexeccond{
		- Existe una versión en estado DRAFT.\\
		- Versión tiene estructura válida o inválida.\\
		- Usuario está autenticado.
		}
		\testcaseexecstep{
		1. Obtener ID de versión.\\
		2. Enviar POST a /api/v1/models/\{model\_id\}/versions/\{version\_id\}/validate con headers.\\
		3. Validar que la respuesta contenga status 200.\\
		4. Respuesta incluye: is\_valid, errors (si los hay), warnings.\\
		5. Si es válida, usar para publicación; si no, mostrar errores.
		}
		\testcaseexpresult{
		- Status HTTP: 200 OK.\\
		- Respuesta incluye estado de validación.\\
		- Listado de errores estructurales/lógicos si existen.\\
		- Listado de advertencias.\\
		- Puede proceder a publicación si is\_valid=true.
		}
		\testcaseuserstory{5}
		\end{acceptancetest}

		\begin{acceptancetest}[tb:HU5_P3]
		\testcasecode{HU5\_P3}
		\testcasename{Publicar versión de modelo}
		\testcasedescription{Validar que una versión validada puede publicarse y pasar a estado PUBLISHED.}
		\testcaseexeccond{
		- Existe versión en estado DRAFT validada.\\
		- Versión no tiene errores críticos.\\
		- Usuario es revisor o administrador.
		}
		\testcaseexecstep{
		1. Obtener ID de versión en estado DRAFT.\\
		2. Enviar POST a /api/v1/models/\{model\_id\}/versions/\{version\_id\}/publish con headers.\\
		3. Validar que la respuesta contenga status 200.\\
		4. Verificar que estado cambia a PUBLISHED.\\
		5. Consultar la versión para confirmar cambio de estado.
		}
		\testcaseexpresult{
		- Status HTTP: 200 OK.\\
		- Estado cambia a PUBLISHED.\\
		- published\_at se establece con timestamp actual.\\
		- Versión anterior (si existe) pasa a estado ARCHIVED (opcional).\\
		- Configuraciones pueden vincularse a esta versión.
		}
		\testcaseuserstory{5}
		\end{acceptancetest}

		\begin{acceptancetest}[tb:HU5_P4]
		\testcasecode{HU5\_P4}
		\testcasename{Recuperar versiones archivadas}
		\testcasedescription{Validar que se pueden listar y recuperar versiones archivadas.}
		\testcaseexeccond{
		- Existen versiones en estado ARCHIVED.\\
		- Usuario está autenticado.\\
		- Usuario tiene permisos para ver historial.
		}
		\testcaseexecstep{
		1. Enviar GET a /api/v1/models/\{model\_id\}/versions?status=ARCHIVED con headers.\\
		2. Validar que la respuesta contenga status 200.\\
		3. Verificar que se retornan solo versiones ARCHIVED.\\
		4. Cada versión incluye: id, version\_number, status, archived\_at, created\_at.\\
		5. Seleccionar una versión y enviar GET a /api/v1/models/\{model\_id\}/versions/\{version\_id\}.
		}
		\testcaseexpresult{
		- Status HTTP: 200 OK.\\
		- Array de versiones archivadas.\\
		- Información completa de cada versión.\\
		- Timestamps de archivación incluidos.\\
		- Puede recuperar estructura de versión archivada.
		}
		\testcaseuserstory{5}
		\end{acceptancetest}

		\subsubsection*{HU-06: Análisis de Métricas y Exportación}

		\begin{acceptancetest}[tb:HU6_P1]
		\testcasecode{HU6\_P1}
		\testcasename{Calcular métricas de versión publicada}
		\testcasedescription{Validar que se pueden calcular métricas de la versión publicada más reciente.}
		\testcaseexeccond{
		- Existe un modelo con versión PUBLISHED.\\
		- La versión tiene estructura completa.\\
		- Usuario está autenticado.
		}
		\testcaseexecstep{
		1. Enviar GET a /api/v1/models/\{model\_id\}/metrics con headers.\\
		2. Validar que la respuesta contenga status 200.\\
		3. Verificar que incluye métricas: feature\_count, group\_count, constraint\_count, cyclomatic\_complexity, tree\_depth.\\
		4. Verificar que se calcula sobre versión PUBLISHED.
		}
		\testcaseexpresult{
		- Status HTTP: 200 OK.\\
		- Respuesta incluye todas las métricas.\\
		- Valores son números enteros >= 0.\\
		- Cálculo se realiza sobre versión publicada.\\
		- Valores coinciden con estructura del modelo.
		}
		\testcaseuserstory{6}
		\end{acceptancetest}

		\begin{acceptancetest}[tb:HU6_P2]
		\testcasecode{HU6\_P2}
		\testcasename{Exportar modelo en formato UVL}
		\testcasedescription{Validar que el modelo se puede exportar en formato UVL válido.}
		\testcaseexeccond{
		- Existe un modelo con versión PUBLISHED.\\
		- La versión tiene estructura completa.\\
		- Usuario está autenticado.
		}
		\testcaseexecstep{
		1. Enviar GET a /api/v1/models/\{model\_id\}/export?format=uvl con headers.\\
		2. Validar que la respuesta contenga status 200.\\
		3. Verificar que el Content-Type es text/plain o application/x-uvl.\\
		4. Validar que el contenido es UVL válido (comienza con features o namespace, tiene sintaxis correcta).\\
		5. Guardar contenido en archivo .uvl y validar.
		}
		\testcaseexpresult{
		- Status HTTP: 200 OK.\\
		- Content-Type correcto.\\
		- UVL contiene sintaxis válida.\\
		- Todas las features, grupos y restricciones están representadas.\\
		- Puede descargarse como archivo.
		}
		\testcaseuserstory{6}
		\end{acceptancetest}

		\begin{acceptancetest}[tb:HU6_P3]
		\testcasecode{HU6\_P3}
		\testcasename{Exportar modelo en formato JSON}
		\testcasedescription{Validar que el modelo se puede exportar en formato JSON.}
		\testcaseexeccond{
		- Existe un modelo con versión PUBLISHED.\\
		- Usuario está autenticado.
		}
		\testcaseexecstep{
		1. Enviar GET a /api/v1/models/\{model\_id\}/export?format=json con headers.\\
		2. Validar que la respuesta contenga status 200.\\
		3. Verificar que Content-Type es application/json.\\
		4. Validar que es JSON válido (se puede parsear).\\
		5. Verificar que incluye: model, version, features, groups, constraints.
		}
		\testcaseexpresult{
		- Status HTTP: 200 OK.\\
		- JSON válido y bien formado.\\
		- Estructura incluye todos los componentes del modelo.\\
		- Puede descargarse como .json.
		}
		\testcaseuserstory{6}
		\end{acceptancetest}

		\begin{acceptancetest}[tb:HU6_P4]
		\testcasecode{HU6\_P4}
		\testcasename{Exportar modelo en formato \ac{DOT}}
		\testcasedescription{Validar que el modelo se puede exportar en formato \ac{DOT} para visualización.}
		\testcaseexeccond{
		- Existe un modelo con versión PUBLISHED.\\
		- Usuario está autenticado.\\
		- El sistema tiene funcionalidad de generación de grafos.
		}
		\testcaseexecstep{
		1. Enviar GET a /api/v1/models/\{model\_id\}/export?format=dot con headers.\\
		2. Validar que la respuesta contenga status 200.\\
		3. Verificar que Content-Type es text/plain o text/vnd.graphviz.\\
		4. Validar sintaxis \ac{DOT} (comienza con digraph o graph).\\
		5. Intentar renderizar con Graphviz para validar.
		}
		\testcaseexpresult{
		- Status HTTP: 200 OK.\\
		- Sintaxis \ac{DOT} válida.\\
		- Nodos representan features.\\
		- Aristas representan relaciones jerárquicas.\\
		- Puede visualizarse con herramientas Graphviz.
		}
		\testcaseuserstory{6}
		\end{acceptancetest}

		\begin{acceptancetest}[tb:HU6_P5]
		\testcasecode{HU6\_P5}
		\testcasename{Exportar modelo en formato Mermaid}
		\testcasedescription{Validar que el modelo se puede exportar en formato Mermaid para diagramas interactivos.}
		\testcaseexeccond{
		- Existe un modelo con versión PUBLISHED.\\
		- Usuario está autenticado.
		}
		\testcaseexecstep{
		1. Enviar GET a /api/v1/models/\{model\_id\}/export?format=mermaid con headers.\\
		2. Validar que la respuesta contenga status 200.\\
		3. Validar sintaxis Mermaid (comienza con graph o flowchart).\\
		4. Verificar que incluye features como nodos.\\
		5. Verificar que relaciones están correctamente representadas.
		}
		\testcaseexpresult{
		- Status HTTP: 200 OK.\\
		- Sintaxis Mermaid válida.\\
		- Diagrama representa jerarquía del modelo.\\
		- Compatible con Mermaid.js.
		}
		\testcaseuserstory{6}
		\end{acceptancetest}

		\subsubsection*{Épica 4: Modelado de Variabilidad Estructural}
		\subsubsection*{HU-07: Edición de la Jerarquía de Características}

		\begin{acceptancetest}[tb:HU7_P1]
		\testcasecode{HU7\_P1}
		\testcasename{Crear feature en versión de modelo}
		\testcasedescription{Validar que se puede crear una nueva feature en una versión específica.}
		\testcaseexeccond{
		- Existe un modelo con versión DRAFT.\\
		- Usuario es diseñador del modelo.\\
		- Usuario está autenticado.
		}
		\testcaseexecstep{
		1. Obtener ID de modelo e ID de versión.\\
		2. Enviar POST a /api/v1/models/\{model\_id\}/versions/\{version\_id\}/features con payload: \{"name": "Módulo A", "type": "CONCRETE"\}.\\
		3. Validar que la respuesta contenga status 201.\\
		4. Verificar que se retorna id, name, type, version\_id, parent\_id.\\
		5. Validar que existe en BD.
		}
		\testcaseexpresult{
		- Status HTTP: 201 Created.\\
		- Feature creada con ID generado.\\
		- Type es CONCRETE, ABSTRACT u OPTIONAL según especificación.\\
		- Sin padre (raíz) inicialmente.\\
		- Timestamps se establecen.
		}
		\testcaseuserstory{7}
		\end{acceptancetest}

		\begin{acceptancetest}[tb:HU7_P2]
		\testcasecode{HU7\_P2}
		\testcasename{Reubicar feature cambiando padre}
		\testcasedescription{Validar que una feature puede ser reubicada cambiando su padre.}
		\testcaseexeccond{
		- Existe una feature con padre.\\
		- Existe otra feature que será el nuevo padre.\\
		- Cambio no crearía ciclo.\\
		- Usuario es diseñador.\\
		- Usuario está autenticado.
		}
		\testcaseexecstep{
		1. Obtener ID de feature a mover y ID de nuevo padre.\\
		2. Enviar PUT a /api/v1/models/\{model\_id\}/versions/\{version\_id\}/features/\{feature\_id\} con payload: \{"parent\_id": "<new\_parent\_id>"\}.\\
		3. Validar que la respuesta contenga status 200.\\
		4. Consultar feature para confirmar nuevo padre.\\
		5. Verificar que subárbol se movió correctamente.
		}
		\testcaseexpresult{
		- Status HTTP: 200 OK.\\
		- Padre actualizado.\\
		- Subárbol completo se mueve con la feature.\\
		- Estructura jerarquía mantiene integridad.\\
		- No se crea ciclo.
		}
		\testcaseuserstory{7}
		\end{acceptancetest}

		\begin{acceptancetest}[tb:HU7_P3]
		\testcasecode{HU7\_P3}
		\testcasename{Recuperar subárbol completo de feature}
		\testcasedescription{Validar que se puede obtener una feature con toda su estructura descendente.}
		\testcaseexeccond{
		- Existe una feature con descendientes.\\
		- Usuario está autenticado.
		}
		\testcaseexecstep{
		1. Obtener ID de feature.\\
		2. Enviar GET a /api/v1/models/\{model\_id\}/versions/\{version\_id\}/features/\{feature\_id\}?include\_subtree=true con headers.\\
		3. Validar que la respuesta contenga status 200.\\
		4. Verificar que incluye la feature y su array de children.\\
		5. Verificar que cada child incluye sus propios children (recursivo).
		}
		\testcaseexpresult{
		- Status HTTP: 200 OK.\\
		- Feature principal incluida.\\
		- Array children contiene todas las features descendentes.\\
		- Estructura jerárquica completa preservada.\\
		- Profundidad del árbol representada correctamente.
		}
		\testcaseuserstory{7}
		\end{acceptancetest}

		\begin{acceptancetest}[tb:HU7_P4]
		\testcasecode{HU7\_P4}
		\testcasename{Desactivar feature con soporte de versionado}
		\testcasedescription{Validar que una feature puede ser desactivada lógicamente manteniéndola en la BD.}
		\testcaseexeccond{
		- Existe una feature activa.\\
		- Feature no es raíz del árbol.\\
		- Usuario es diseñador.\\
		- Usuario está autenticado.
		}
		\testcaseexecstep{
		1. Obtener ID de feature a desactivar.\\
		2. Enviar PUT a /api/v1/models/\{model\_id\}/versions/\{version\_id\}/features/\{feature\_id\} con payload: \{"is\_active": false\}.\\
		3. Validar que la respuesta contenga status 200.\\
		4. Consultar feature para confirmar is\_active=False.\\
		5. Verificar que no afecta otras versiones.
		}
		\testcaseexpresult{
		- Status HTTP: 200 OK.\\
		- Feature marcada como is\_active=False.\\
		- Feature permanece en BD.\\
		- Otros cambios no afectados.\\
		- Versionado mantiene registro del cambio.
		}
		\testcaseuserstory{7}
		\end{acceptancetest}

		\begin{acceptancetest}[tb:HU7_P5]
		\testcasecode{HU7\_P5}
		\testcasename{Actualizar propiedades de feature}
		\testcasedescription{Validar que se pueden actualizar propiedades de una feature.}
		\testcaseexeccond{
		- Existe una feature.\\
		- Usuario es diseñador.\\
		- Usuario está autenticado.
		}
		\testcaseexecstep{
		1. Obtener ID de feature.\\
		2. Enviar PUT a /api/v1/models/\{model\_id\}/versions/\{version\_id\}/features/\{feature\_id\} con payload: \{"name": "Nuevo nombre", "type": "ABSTRACT"\}.\\
		3. Validar que la respuesta contenga status 200.\\
		4. Consultar feature para confirmar cambios.\\
		5. Verificar que metadata se actualizó.
		}
		\testcaseexpresult{
		- Status HTTP: 200 OK.\\
		- Nombre actualizado.\\
		- Tipo actualizado.\\
		- updated\_at se actualiza.\\
		- Cambios persistidos en BD.
		}
		\testcaseuserstory{7}
		\end{acceptancetest}

		\subsubsection*{HU-08: Modelado de Grupos y Relaciones Estructurales}

		\begin{acceptancetest}[tb:HU8_P1]
		\testcasecode{HU8\_P1}
		\testcasename{Crear grupo XOR entre features}
		\testcasedescription{Validar que se puede crear un grupo XOR (exactamente uno) entre features hermanas.}
		\testcaseexeccond{
		- Existen features hermanas en versión DRAFT.\\
		- Features tienen mismo padre.\\
		- Usuario es diseñador.\\
		- Usuario está autenticado.
		}
		\testcaseexecstep{
		1. Obtener IDs de features a agrupar.\\
		2. Enviar POST a /api/v1/models/\{model\_id\}/versions/\{version\_id\}/groups con payload: \{"group\_type": "XOR", "parent\_id": "<parent\_id>"\}.\\
		3. Validar que la respuesta contenga status 201.\\
		4. Verificar que se retorna id, group\_type, parent\_id, min\_cardinality=1, max\_cardinality=1.\\
		5. Asignar features al grupo.
		}
		\testcaseexpresult{
		- Status HTTP: 201 Created.\\
		- Grupo creado con group\_type=XOR.\\
		- Cardinalidades establecidas: min=1, max=1.\\
		- ID generado.\\
		- Grupo listo para asignar features.
		}
		\testcaseuserstory{8}
		\end{acceptancetest}

		\begin{acceptancetest}[tb:HU8_P2]
		\testcasecode{HU8\_P2}
		\testcasename{Crear grupo OR con cardinalidades}
		\testcasedescription{Validar que se puede crear un grupo OR con cardinalidades mínimas y máximas personalizadas.}
		\testcaseexeccond{
		- Existen features hermanas en versión DRAFT.\\
		- Usuario es diseñador.\\
		- Usuario está autenticado.
		}
		\testcaseexecstep{
		1. Obtener ID de padre.\\
		2. Enviar POST a /api/v1/models/\{model\_id\}/versions/\{version\_id\}/groups con payload: \{"group\_type": "OR", "parent\_id": "<parent\_id>", "min\_cardinality": 1, "max\_cardinality": 3\}.\\
		3. Validar que la respuesta contenga status 201.\\
		4. Verificar que cardinalidades se establecen correctamente.\\
		5. Asignar features al grupo.
		}
		\testcaseexpresult{
		- Status HTTP: 201 Created.\\
		- Grupo creado con group\_type=OR.\\
		- Cardinalidades: min=1, max=3.\\
		- Validación: min >= 1 y max >= min.\\
		- Grupo aceptará entre 1 y 3 features seleccionadas.
		}
		\testcaseuserstory{8}
		\end{acceptancetest}

		\begin{acceptancetest}[tb:HU8_P3]
		\testcasecode{HU8\_P3}
		\testcasename{Crear relación requires entre features}
		\testcasedescription{Validar que se puede crear una relación de implicación (requires) entre features.}
		\testcaseexeccond{
		- Existen dos features en versión DRAFT.\\
		- Features pueden estar en diferentes ramas.\\
		- Usuario es diseñador.\\
		- Usuario está autenticado.
		}
		\testcaseexecstep{
		1. Obtener IDs de feature origen y destino.\\
		2. Enviar POST a /api/v1/models/\{model\_id\}/versions/\{version\_id\}/relationships con payload: \{"relationship\_type": "REQUIRES", "source\_id": "<feature\_id\_1>", "target\_id": "<feature\_id\_2>"\}.\\
		3. Validar que la respuesta contenga status 201.\\
		4. Verificar que se retorna relación con IDs correctos.\\
		5. Guardar ID de relación.
		}
		\testcaseexpresult{
		- Status HTTP: 201 Created.\\
		- Relación creada.\\
		- Tipo: REQUIRES (implicación).\\
		- Source y target correctos.\\
		- Bidireccional en términos de consulta.
		}
		\testcaseuserstory{8}
		\end{acceptancetest}

		\begin{acceptancetest}[tb:HU8_P4]
		\testcasecode{HU8\_P4}
		\testcasename{Crear relación excludes entre features}
		\testcasedescription{Validar que se puede crear una relación de exclusión entre features.}
		\testcaseexeccond{
		- Existen dos features en versión DRAFT.\\
		- Usuario es diseñador.\\
		- Usuario está autenticado.
		}
		\testcaseexecstep{
		1. Obtener IDs de features a relacionar.\\
		2. Enviar POST a /api/v1/models/\{model\_id\}/versions/\{version\_id\}/relationships con payload: \{"relationship\_type": "EXCLUDES", "source\_id": "<id\_1>", "target\_id": "<id\_2>"\}.\\
		3. Validar que la respuesta contenga status 201.\\
		4. Verificar que relación se crea.\\
		5. Intentar crear configuración con ambas features (debe fallar).
		}
		\testcaseexpresult{
		- Status HTTP: 201 Created.\\
		- Relación de exclusión creada.\\
		- Ambas features no pueden seleccionarse simultáneamente.\\
		- Validación de configuraciones lo verifica.
		}
		\testcaseuserstory{8}
		\end{acceptancetest}

		\begin{acceptancetest}[tb:HU8_P5]
		\testcasecode{HU8\_P5}
		\testcasename{Listar grupos y relaciones de versión}
		\testcasedescription{Validar que se pueden obtener todos los grupos y relaciones de una versión.}
		\testcaseexeccond{
		- Existe versión con grupos y relaciones.\\
		- Usuario está autenticado.
		}
		\testcaseexecstep{
		1. Enviar GET a /api/v1/models/\{model\_id\}/versions/\{version\_id\}/groups con headers.\\
		2. Validar que la respuesta contenga status 200.\\
		3. Verificar array de grupos con id, group\_type, parent\_id, min\_cardinality, max\_cardinality.\\
		4. Enviar GET a /api/v1/models/\{model\_id\}/versions/\{version\_id\}/relationships.\\
		5. Verificar array de relaciones.
		}
		\testcaseexpresult{
		- Status HTTP: 200 OK (ambas).\\
		- Todos los grupos listados.\\
		- Todas las relaciones listadas.\\
		- Información completa de cada grupo/relación.
		}
		\testcaseuserstory{8}
		\end{acceptancetest}

		\subsubsection*{Épica 5: Restricciones y Lógica Avanzada}
		\subsubsection*{HU-09: Gestión de Restricciones Transversales}

		\begin{acceptancetest}[tb:HU9_P1]
		\testcasecode{HU9\_P1}
		\testcasename{Crear restricción simples con sintaxis UVL}
		\testcasedescription{Validar que se puede crear una restricción simple con expresión lógica.}
		\testcaseexeccond{
		- Existe versión DRAFT con features.\\
		- Usuario es diseñador.\\
		- Usuario está autenticado.
		}
		\testcaseexecstep{
		1. Obtener ID de versión.\\
		2. Enviar POST a /api/v1/models/\{model\_id\}/versions/\{version\_id\}/constraints con payload: \{"expression": "FeatureA => FeatureB", "description": "Si A entonces B"\}.\\
		3. Validar que la respuesta contenga status 201.\\
		4. Verificar que se parsea y valida la expresión.\\
		5. Guardar ID de restricción.
		}
		\testcaseexpresult{
		- Status HTTP: 201 Created.\\
		- Restricción creada con ID.\\
		- Expresión validada sintácticamente.\\
		- Descripción opcional almacenada.\\
		- Restricción activa.
		}
		\testcaseuserstory{9}
		\end{acceptancetest}

		\begin{acceptancetest}[tb:HU9_P2]
		\testcasecode{HU9\_P2}
		\testcasename{Crear restricción compleja con múltiples operadores}
		\testcasedescription{Validar que se pueden crear restricciones complejas con operadores lógicos.}
		\testcaseexeccond{
		- Existe versión DRAFT con features.\\
		- Usuario es diseñador.\\
		- Usuario está autenticado.
		}
		\testcaseexecstep{
		1. Obtener ID de versión.\\
		2. Enviar POST a /api/v1/models/\{model\_id\}/versions/\{version\_id\}/constraints con payload: \{"expression": "(FeatureA | FeatureB) \& \textasciitilde FeatureC => FeatureD", "description": "Restricción compleja"\}.\\
		3. Validar que la respuesta contenga status 201.\\
		4. Verificar que se valida correctamente.\\
		5. Intentar añadir a validación del modelo.
		}
		\testcaseexpresult{
		- Status HTTP: 201 Created.\\
		- Restricción con múltiples operadores aceptada.\\
		- Sintaxis: |, \&, \textasciitilde, =>, válidas.\\
		- Expresión se almacena correctamente.
		}
		\testcaseuserstory{9}
		\end{acceptancetest}

		\begin{acceptancetest}[tb:HU9_P3]
		\testcasecode{HU9\_P3}
		\testcasename{Actualizar restricción existente}
		\testcasedescription{Validar que una restricción puede ser actualizada con nueva expresión.}
		\testcaseexeccond{
		- Existe restricción en versión DRAFT.\\
		- Usuario es diseñador.\\
		- Usuario está autenticado.
		}
		\testcaseexecstep{
		1. Obtener ID de restricción.\\
		2. Enviar PUT a /api/v1/models/\{model\_id\}/versions/\{version\_id\}/constraints/\{constraint\_id\} con payload: \{"expression": "FeatureA \& FeatureB => FeatureC", "description": "Expresión actualizada"\}.\\
		3. Validar que la respuesta contenga status 200.\\
		4. Consultar restricción para confirmar cambios.\\
		5. Verificar que versionado registra el cambio.
		}
		\testcaseexpresult{
		- Status HTTP: 200 OK.\\
		- Expresión actualizada.\\
		- Nueva expresión validada.\\
		- Descripción actualizada.\\
		- updated\_at se modifica.
		}
		\testcaseuserstory{9}
		\end{acceptancetest}

		\begin{acceptancetest}[tb:HU9_P4]
		\testcasecode{HU9\_P4}
		\testcasename{Listar restricciones de versión}
		\testcasedescription{Validar que se pueden obtener todas las restricciones de una versión.}
		\testcaseexeccond{
		- Existe versión con restricciones.\\
		- Usuario está autenticado.
		}
		\testcaseexecstep{
		1. Enviar GET a /api/v1/models/\{model\_id\}/versions/\{version\_id\}/constraints con headers.\\
		2. Validar que la respuesta contenga status 200.\\
		3. Verificar array de restricciones.\\
		4. Cada restricción incluye: id, expression, description, is\_active, created\_at.
		}
		\testcaseexpresult{
		- Status HTTP: 200 OK.\\
		- Array de todas las restricciones.\\
		- Información completa de cada restricción.\\
		- Ordenadas por fecha de creación.
		}
		\testcaseuserstory{9}
		\end{acceptancetest}

		\subsubsection*{HU-010: Integración y Sincronización \ac{UVL}}
		
		\begin{acceptancetest}[tb:HU10_P1]
		\testcasecode{HU10\_P1}
		\testcasename{Aplicar contenido \ac{UVL} a estructura de modelo}
		\testcasedescription{Validar que se puede cargar un modelo completo desde contenido \ac{UVL}.}
		\testcaseexeccond{
		- Existe una versión DRAFT vacía.\\
		- Se dispone de contenido \ac{UVL} válido.\\
		- Usuario es diseñador.\\
		- Usuario está autenticado.
		}
		\testcaseexecstep{
		1. Preparar contenido \ac{UVL} válido (features, grupos, restricciones).\\
		2. Enviar PUT a /api/v1/models/\{model\_id\}/versions/\{version\_id\}/uvl con payload: \{"content": "<contenido UVL>"\}.\\
		3. Validar que la respuesta contenga status 200.\\
		4. Consultar la estructura generada del modelo.\\
		5. Verificar que árbol, grupos y restricciones se crearon correctamente.
		}
		\testcaseexpresult{
		- Status HTTP: 200 OK.\\
		- Estructura del modelo se genera desde \ac{UVL}.\\
		- Todas las features se crean.\\
		- Todos los grupos se crean.\\
		- Todas las restricciones se crean.\\
		- Estructura coincide con \ac{UVL} original.
		}
		\testcaseuserstory{10}
		\end{acceptancetest}

		\begin{acceptancetest}[tb:HU10_P2]
		\testcasecode{HU10\_P2}
		\testcasename{Exportar contenido \ac{UVL} desde modelo}
		\testcasedescription{Validar que se puede generar contenido \ac{UVL} válido desde la estructura del modelo.}
		\testcaseexeccond{
		- Existe versión PUBLISHED con estructura completa.\\
		- Estructura es válida.\\
		- Usuario está autenticado.
		}
		\testcaseexecstep{
		1. Enviar GET a /api/v1/models/\{model\_id\}/export?format=uvl con headers.\\
		2. Validar que la respuesta contenga status 200.\\
		3. Extraer contenido \ac{UVL}.\\
		4. Validar sintaxis \ac{UVL}.\\
		5. Importar \ac{UVL} en nueva versión y comparar estructura.
		}
		\testcaseexpresult{
		- Status HTTP: 200 OK.\\
		- UVL válido y bien formado.\\
		- Contiene todas las features, grupos y restricciones.\\
		- Re-importación genera estructura idéntica.
		}
		\testcaseuserstory{10}
		\end{acceptancetest}

		\begin{acceptancetest}[tb:HU10_P3]
		\testcasecode{HU10\_P3}
		\testcasename{Persistir UVL analizado en versión}
		\testcasedescription{Validar que el contenido UVL se analiza, valida y persiste en la BD.}
		\testcaseexeccond{
		- Existe versión DRAFT.\\
		- Se carga contenido UVL válido.\\
		- Usuario es diseñador.\\
		- Usuario está autenticado.
		}
		\testcaseexecstep{
		1. Enviar PUT a /api/v1/models/\{model\_id\}/versions/\{version\_id\}/uvl con contenido UVL.\\
		2. Validar que la respuesta contenga status 200.\\
		3. Consultar GET /api/v1/models/\{model\_id\}/versions/\{version\_id\}/uvl.\\
		4. Verificar que el UVL almacenado coincide con el original.\\
		5. Validar que existe en BD.
		}
		\testcaseexpresult{
		- Status HTTP: 200 OK (ambas).\\
		- UVL almacenado en BD.\\
		- Análisis y tokens se persisten.\\
		- Puede recuperarse completo.\\
		- Válido para re-importación.
		}
		\testcaseuserstory{10}
		\end{acceptancetest}

		\subsubsection*{Épica 6: Validación y Análisis de Modelos}
		\subsubsection*{HU-11: Validación Lógica y Estructural del Modelo}
		
		\begin{acceptancetest}[tb:HU11_P1]
		\testcasecode{HU11\_P1}
		\testcasename{Evaluar satisfacibilidad global del modelo}
		\testcasedescription{Validar que se puede ejecutar un análisis SAT sobre el modelo completo.}
		\testcaseexeccond{
		- Existe versión con estructura y restricciones.\\
		- Sistema tiene integración con solucionador SAT.\\
		- Usuario está autenticado.
		}
		\testcaseexecstep{
		1. Enviar POST a /api/v1/models/\{model\_id\}/versions/\{version\_id\}/validate/satisfiability con headers.\\
		2. Validar que la respuesta contenga status 200.\\
		3. Verificar que respuesta incluye: is\_satisfiable, error\_message (si aplica).\\
		4. Si es satisfiable, debe existir al menos una configuración válida.\\
		5. Si no es satisfiable, revisar restricciones.
		}
		\testcaseexpresult{
		- Status HTTP: 200 OK.\\
		- Respuesta incluye is\_satisfiable (true/false).\\
		- Si false, mensaje de error explicando por qué.\\
		- Cálculo se ejecuta sobre modelo completo.\\
		- Tiempo de ejecución razonable.
		}
		\testcaseuserstory{11}
		\end{acceptancetest}

		\begin{acceptancetest}[tb:HU11_P2]
		\testcasecode{HU11\_P2}
		\testcasename{Detectar ciclos en la jerarquía}
		\testcasedescription{Validar que se detectan ciclos en la relaciones entre features.}
		\testcaseexeccond{
		- Existe versión con estructura.\\
		- Estructura puede contener ciclos (intencional para testing).\\
		- Usuario está autenticado.
		}
		\testcaseexecstep{
		1. Crear escenario con ciclo: A->B->C->A.\\
		2. Enviar POST a /api/v1/models/\{model\_id\}/versions/\{version\_id\}/validate/structure con headers.\\
		3. Validar que la respuesta contenga status 200.\\
		4. Verificar que se reporta error: "Cycle detected: A->B->C->A".\\
		5. Listar features participantes.
		}
		\testcaseexpresult{
		- Status HTTP: 200 OK.\\
		- Respuesta incluye lista de ciclos detectados.\\
		- Cada ciclo lista features que participan.\\
		- Mensaje claro de error.
		}
		\testcaseuserstory{11}
		\end{acceptancetest}

		\begin{acceptancetest}[tb:HU11_P3]
		\testcasecode{HU11\_P3}
		\testcasename{Identificar características muertas}
		\testcasedescription{Validar que se detectan características que nunca pueden ser seleccionadas (dead features).}
		\testcaseexeccond{
		- Existe versión con restricciones que hacen imposible seleccionar cierta feature.\\
		- Usuario está autenticado.
		}
		\testcaseexecstep{
		1. Crear estructura con feature imposible de seleccionar.\\
		2. Enviar POST a /api/v1/models/\{model\_id\}/versions/\{version\_id\}/analyze/dead-features con headers.\\
		3. Validar que la respuesta contenga status 200.\\
		4. Verificar array de features muertas.\\
		5. Cada feature incluye razón por la que es muerta.
		}
		\testcaseexpresult{
		- Status HTTP: 200 OK.\\
		- Array de features muertas identificadas.\\
		- Explicación de por qué es muerta.\\
		- Análisis completo y correcto.
		}
		\testcaseuserstory{11}
		\end{acceptancetest}

		\begin{acceptancetest}[tb:HU11_P4]
		\testcasecode{HU11\_P4}
		\testcasename{Detectar nodos huérfanos}
		\testcasedescription{Validar que se detectan features que no tienen ruta a raíz.}
		\testcaseexeccond{
		- Existe versión con estructura.\\
		- Se intenta crear feature sin padre (orphan).\\
		- Usuario está autenticado.
		}
		\testcaseexecstep{
		1. Buscar features sin ruta a raíz.\\
		2. Enviar POST a /api/v1/models/\{model\_id\}/versions/\{version\_id\}/validate/structure con headers.\\
		3. Validar que la respuesta reporta nodos huérfanos.\\
		4. Verificar que lista los IDs de nodos huérfanos.\\
		5. Mensaje explicativo.
		}
		\testcaseexpresult{
		- Status HTTP: 200 OK.\\
		- Lista de nodos huérfanos detectados.\\
		- Explicación clara.\\
		- Recomendación de reparación.
		}
		\testcaseuserstory{11}
		\end{acceptancetest}

		\begin{acceptancetest}[tb:HU11_P5]
		\testcasecode{HU11\_P5}
		\testcasename{Ejecutar validación completa}
		\testcasedescription{Validar que se ejecuta suite completa de validaciones en un paso.}
		\testcaseexeccond{
		- Existe versión con estructura.\\
		- Usuario está autenticado.
		}
		\testcaseexecstep{
		1. Enviar POST a /api/v1/models/\{model\_id\}/versions/\{version\_id\}/validate/full con headers.\\
		2. Validar que la respuesta contenga status 200.\\
		3. Respuesta incluye: satisfiable, cycles, dead\_features, orphan\_nodes, warnings.\\
		4. Cada sección tiene lista de errores/warnings.\\
		5. Resumen general de validación.
		}
		\testcaseexpresult{
		- Status HTTP: 200 OK.\\
		- Respuesta incluye todas las validaciones.\\
		- Formato consistente.\\
		- Fácil de interpretar.\\
		- Guía para correcciones.
		}
		\testcaseuserstory{11}
		\end{acceptancetest}

		\subsubsection*{Épica 7: Configuración, Optimización y Recursos}
		\subsubsection*{HU-12: Gestión de Configuraciones}

		\begin{acceptancetest}[tb:HU12_P1]
		\testcasecode{HU12\_P1}
		\testcasename{Crear configuración en versión publicada}
		\testcasedescription{Validar que se puede crear una configuración (selección de features) en una versión publicada.}
		\testcaseexeccond{
		- Existe versión PUBLISHED con features.\\
		- Se dispone de lista válida de features.\\
		- Usuario es configurador.\\
		- Usuario está autenticado.
		}
		\testcaseexecstep{
		1. Obtener IDs de features válidas.\\
		2. Enviar POST a /api/v1/configurations con payload: \{"name": "Config 1", "description": "Configuración de prueba", "version\_id": "<version\_id>", "features": ["<id\_1>", "<id\_2>"]\}.\\
		3. Validar que la respuesta contenga status 201.\\
		4. Verificar que se crea configuración con ID.\\
		5. Validar que se vincula a versión PUBLISHED.
		}
		\testcaseexpresult{
		- Status HTTP: 201 Created.\\
		- Configuración creada con ID.\\
		- Vinculada a versión PUBLISHED.\\
		- Selección de features persistida.\\
		- Timestamps se establecen.
		}
		\testcaseuserstory{12}
		\end{acceptancetest}

		\begin{acceptancetest}[tb:HU12_P2]
		\testcasecode{HU12\_P2}
		\testcasename{Actualizar configuración existente}
		\testcasedescription{Validar que se puede actualizar nombre, descripción y features de una configuración.}
		\testcaseexeccond{
		- Existe configuración persistida.\\
		- Usuario es propietario o administrador.\\
		- Usuario está autenticado.
		}
		\testcaseexecstep{
		1. Obtener ID de configuración.\\
		2. Enviar PUT a /api/v1/configurations/\{config\_id\} con payload: \{"name": "Nuevo nombre", "description": "Descripción actualizada", "features": ["<nuevo\_id\_1>", "<nuevo\_id\_2>"]\}.\\
		3. Validar que la respuesta contenga status 200.\\
		4. Consultar configuración para confirmar cambios.\\
		5. Verificar que features se actualizaron.
		}
		\testcaseexpresult{
		- Status HTTP: 200 OK.\\
		- Configuración actualizada en BD.\\
		- Nombre y descripción cambiados.\\
		- Lista de features actualizada.\\
		- updated\_at se modifica.
		}
		\testcaseuserstory{12}
		\end{acceptancetest}

		\begin{acceptancetest}[tb:HU12_P3]
		\testcasecode{HU12\_P3}
		\testcasename{Listar configuraciones paginadas}
		\testcasedescription{Validar que se pueden obtener configuraciones con paginación.}
		\testcaseexeccond{
		- Existen múltiples configuraciones persistidas.\\
		- Usuario está autenticado.
		}
		\testcaseexecstep{
		1. Enviar GET a /api/v1/configurations?skip=0\&limit=10 con headers.\\
		2. Validar que la respuesta contenga status 200.\\
		3. Verificar array de configuraciones.\\
		4. Cada config incluye: id, name, description, version\_id, created\_at, updated\_at.\\
		5. Validar paginación: total, skip, limit.
		}
		\testcaseexpresult{
		- Status HTTP: 200 OK.\\
		- Array de configuraciones paginadas.\\
		- Información completa de cada configuración.\\
		- Paginación funciona correctamente.\\
		- Total de registros incluido.
		}
		\testcaseuserstory{12}
		\end{acceptancetest}

		\begin{acceptancetest}[tb:HU12_P4]
		\testcasecode{HU12\_P4}
		\testcasename{Filtrar configuraciones por versión}
		\testcasedescription{Validar que se pueden filtrar configuraciones por ID de versión.}
		\testcaseexeccond{
		- Existen configuraciones de múltiples versiones.\\
		- Usuario está autenticado.
		}
		\testcaseexecstep{
		1. Obtener ID de versión específica.\\
		2. Enviar GET a /api/v1/configurations?version\_id=<version\_id>\&skip=0\&limit=10 con headers.\\
		3. Validar que la respuesta contenga status 200.\\
		4. Verificar que todas las configuraciones retornadas tienen version\_id correcto.\\
		5. Validar paginación.
		}
		\testcaseexpresult{
		- Status HTTP: 200 OK.\\
		- Solo configuraciones de versión especificada.\\
		- Paginación correcta.\\
		- Filtrado preciso.
		}
		\testcaseuserstory{12}
		\end{acceptancetest}

		\begin{acceptancetest}[tb:HU12_P5]
		\testcasecode{HU12\_P5}
		\testcasename{Desactivar configuración}
		\testcasedescription{Validar que se puede desactivar una configuración sin eliminarla.}
		\testcaseexeccond{
		- Existe configuración activa.\\
		- Usuario es propietario o administrador.\\
		- Usuario está autenticado.
		}
		\testcaseexecstep{
		1. Obtener ID de configuración.\\
		2. Enviar PUT a /api/v1/configurations/\{config\_id\} con payload: \{"is\_active": false\}.\\
		3. Validar que la respuesta contenga status 200.\\
		4. Consultar configuración para confirmar is\_active=False.\\
		5. Verificar que datos persisten en BD.
		}
		\testcaseexpresult{
		- Status HTTP: 200 OK.\\
		- Configuración marcada como is\_active=False.\\
		- Datos persistidos.\\
		- No aparece en listados de activas (opcional).
		}
		\testcaseuserstory{12}
		\end{acceptancetest}

		\subsubsection*{HU-13: Validación y Generación Automática de Configuraciones}

		\begin{acceptancetest}[tb:HU13_P1]
		\testcasecode{HU13\_P1}
		\testcasename{Validar selección sin persistirla}
		\testcasedescription{Validar que se puede verificar si una selección es consistente sin guardarla.}
		\testcaseexeccond{
		- Existe versión PUBLISHED con estructura y restricciones.\\
		- Se dispone de lista de features.\\
		- Usuario está autenticado.
		}
		\testcaseexecstep{
		1. Enviar POST a /api/v1/configurations/validate con payload: \{"version\_id": "<version\_id>", "features": ["<id\_1>", "<id\_2>"]\}.\\
		2. Validar que la respuesta contenga status 200.\\
		3. Respuesta incluye: is\_valid, errors (si existen), warnings.\\
		4. Verificar que no se crea configuración persistida.\\
		5. Mensaje claro de validación.
		}
		\testcaseexpresult{
		- Status HTTP: 200 OK.\\
		- Respuesta indica validez.\\
		- Lista de errores si los hay.\\
		- Lista de advertencias.\\
		- No se persiste la configuración.
		}
		\testcaseuserstory{13}
		\end{acceptancetest}

		\begin{acceptancetest}[tb:HU13_P2]
		\testcasecode{HU13\_P2}
		\testcasename{Generar configuraciones válidas automáticamente}
		\testcasedescription{Validar que el sistema puede generar automáticamente configuraciones válidas.}
		\testcaseexeccond{
		- Existe versión PUBLISHED satisfiable.\\
		- Usuario está autenticado.\\
		- Sistema tiene solucionador SAT integrado.
		}
		\testcaseexecstep{
		1. Enviar POST a /api/v1/configurations/generate con payload: \{"version\_id": "<version\_id>", "limit": 5\}.\\
		2. Validar que la respuesta contenga status 200.\\
		3. Respuesta incluye array de configuraciones válidas.\\
		4. Cada una con lista de features seleccionadas.\\
		5. Límite de generaciones respetado.
		}
		\testcaseexpresult{
		- Status HTTP: 200 OK.\\
		- Array de configuraciones generadas.\\
		- Cada una es válida según restricciones.\\
		- Cantidad <= limit.\\
		- Diversidad de soluciones (si limit > 1).
		}
		\testcaseuserstory{13}
		\end{acceptancetest}

		\begin{acceptancetest}[tb:HU13_P3]
		\testcasecode{HU13\_P3}
		\testcasename{Generar configuración óptima}
		\testcasedescription{Validar que se puede generar una configuración optimizada según criterios.}
		\testcaseexeccond{
		- Existe versión PUBLISHED.\\
		- Se definen criterios de optimización.\\
		- Usuario está autenticado.
		}
		\testcaseexecstep{
		1. Enviar POST a /api/v1/configurations/optimize con payload: \{"version\_id": "<version\_id>", "objective": "maximize\_features"\}.\\
		2. Validar que la respuesta contenga status 200.\\
		3. Respuesta incluye configuración optimizada.\\
		4. Incluye métricas de optimización.\\
		5. Configuración es válida.
		}
		\testcaseexpresult{
		- Status HTTP: 200 OK.\\
		- Configuración optimizada devuelta.\\
		- Cumple restricciones.\\
		- Optimización aplicada.\\
		- Métricas incluidas.
		}
		\testcaseuserstory{13}
		\end{acceptancetest}

		\begin{acceptancetest}[tb:HU13_P4]
		\testcasecode{HU13\_P4}
		\testcasename{Validar configuración contra restricciones complejas}
		\testcasedescription{Validar que la validación considera restricciones complejas con múltiples operadores.}
		\testcaseexeccond{
		- Existe versión con restricciones complejas.\\
		- Se dispone de selección a validar.\\
		- Usuario está autenticado.
		}
		\testcaseexecstep{
		1. Crear restricción compleja: (A | B) \& \textasciitilde C => D.\\
		2. Enviar POST a /api/v1/configurations/validate con selección que viola.\\
		3. Validar que la respuesta reporta error.\\
		4. Mensaje indica qué restricción se viola.\\
		5. Sugerir alternativa válida.
		}
		\testcaseexpresult{
		- Status HTTP: 200 OK.\\
		- is\_valid=false.\\
		- Error específico reportado.\\
		- Restricción violada identificada.\\
		- Sugerencia de corrección.
		}
		\testcaseuserstory{13}
		\end{acceptancetest}

		\subsubsection*{Épica 6: Validación y Análisis de Modelos}
		\subsubsection*{HU-11: Enriquecimiento con Recursos y Taxonomías}

		\begin{acceptancetest}[tb:HU14_P1]
		\testcasecode{HU14\_P1}
		\testcasename{Crear etiqueta (tag)}
		\testcasedescription{Validar que se pueden crear etiquetas para clasificar features.}
		\testcaseexeccond{
		- Existe modelo.\\
		- Usuario es diseñador o administrador.\\
		- Usuario está autenticado.
		}
		\testcaseexecstep{
		1. Enviar POST a /api/v1/tags con payload: \{"name": "Programación", "color": "\#FF5733"\}.\\
		2. Validar que la respuesta contenga status 201.\\
		3. Verificar que se retorna ID de etiqueta.\\
		4. Verificar que se persiste en BD.\\
		5. Guardar ID para asociar a features.
		}
		\testcaseexpresult{
		- Status HTTP: 201 Created.\\
		- Etiqueta creada con ID.\\
		- Nombre y color se almacenan.\\
		- Puede utilizarse para asociaciones.
		}
		\testcaseuserstory{14}
		\end{acceptancetest}

		\begin{acceptancetest}[tb:HU14_P2]
		\testcasecode{HU14\_P2}
		\testcasename{Asociar etiqueta a feature}
		\testcasedescription{Validar que se puede asociar una etiqueta a una feature.}
		\testcaseexeccond{
		- Existe feature en versión.\\
		- Existe etiqueta creada.\\
		- Usuario es diseñador.\\
		- Usuario está autenticado.
		}
		\testcaseexecstep{
		1. Obtener IDs de feature y etiqueta.\\
		2. Enviar POST a /api/v1/models/\{model\_id\}/versions/\{version\_id\}/features/\{feature\_id\}/tags/\{tag\_id\} con headers.\\
		3. Validar que la respuesta contenga status 200.\\
		4. Consultar feature para confirmar asociación.\\
		5. Verificar que tag aparece en lista de tags de feature.
		}
		\testcaseexpresult{
		- Status HTTP: 200 OK.\\
		- Etiqueta asociada a feature.\\
		- Persistida en BD.\\
		- Consultas posteriores la incluyen.
		}
		\testcaseuserstory{14}
		\end{acceptancetest}

		\begin{acceptancetest}[tb:HU14_P3]
		\testcasecode{HU14\_P3}
		\testcasename{Desasociar etiqueta de feature}
		\testcasedescription{Validar que se puede remover una etiqueta asociada a una feature.}
		\testcaseexeccond{
		- Feature tiene etiquetas asociadas.\\
		- Usuario es diseñador.\\
		- Usuario está autenticado.
		}
		\testcaseexecstep{
		1. Obtener IDs de feature y tag asociado.\\
		2. Enviar DELETE a /api/v1/models/\{model\_id\}/versions/\{version\_id\}/features/\{feature\_id\}/tags/\{tag\_id\} con headers.\\
		3. Validar que la respuesta contenga status 200.\\
		4. Consultar feature para confirmar que tag no está asociado.\\
		5. Verificar que se removió de BD.
		}
		\testcaseexpresult{
		- Status HTTP: 200 OK.\\
		- Etiqueta removida de feature.\\
		- BD actualizada.\\
		- Tag sigue existiendo (se reutiliza).
		}
		\testcaseuserstory{14}
		\end{acceptancetest}

		\begin{acceptancetest}[tb:HU14_P4]
		\testcasecode{HU14\_P4}
		\testcasename{Crear recurso educativo}
		\testcasedescription{Validar que se puede crear un recurso (archivo, enlace) asociado al modelo.}
		\testcaseexeccond{
		- Existe modelo.\\
		- Usuario es editor de recursos.\\
		- Usuario está autenticado.\\
		- Se dispone de archivo o URL.
		}
		\testcaseexecstep{
		1. Enviar POST a /api/v1/resources con payload multipart: \{name: "Introducción a Java", type: "VIDEO", url: "https://...", language: "es", license: "CC-BY-4.0"\}.\\
		2. Validar que la respuesta contenga status 201.\\
		3. Verificar que se retorna ID de recurso.\\
		4. Verificar metadatos almacenados.\\
		5. Guardar ID para asociar a features/configuraciones.
		}
		\testcaseexpresult{
		- Status HTTP: 201 Created.\\
		- Recurso creado con ID.\\
		- Metadatos almacenados.\\
		- Propietario registrado automáticamente.\\
		- Timestamps se establecen.
		}
		\testcaseuserstory{14}
		\end{acceptancetest}

		\begin{acceptancetest}[tb:HU14_P5]
		\testcasecode{HU14\_P5}
		\testcasename{Asociar recurso a feature}
		\testcasedescription{Validar que se puede vincular un recurso educativo a una feature.}
		\testcaseexeccond{
		- Existe feature en versión.\\
		- Existe recurso creado.\\
		- Usuario es editor.\\
		- Usuario está autenticado.
		}
		\testcaseexecstep{
		1. Obtener IDs de feature y recurso.\\
		2. Enviar POST a /api/v1/models/\{model\_id\}/versions/\{version\_id\}/features/\{feature\_id\}/resources/\{resource\_id\} con headers.\\
		3. Validar que la respuesta contenga status 200.\\
		4. Consultar feature para verificar asociación.\\
		5. Recurso aparece en lista de recursos de feature.
		}
		\testcaseexpresult{
		- Status HTTP: 200 OK.\\
		- Recurso asociado a feature.\\
		- Persistencia en BD.\\
		- Versiones se registran.
		}
		\testcaseuserstory{14}
		\end{acceptancetest}

		\begin{acceptancetest}[tb:HU14_P6]
		\testcasecode{HU14\_P6}
		\testcasename{Actualizar metadatos de recurso}
		\testcasedescription{Validar que se pueden actualizar metadatos del recurso (idioma, licencia, etc).}
		\testcaseexeccond{
		- Existe recurso persistido.\\
		- Usuario es propietario del recurso.\\
		- Usuario está autenticado.
		}
		\testcaseexecstep{
		1. Obtener ID de recurso.\\
		2. Enviar PUT a /api/v1/resources/\{resource\_id\} con payload: \{"language": "en", "license": "MIT", "description": "Nueva descripción"\}.\\
		3. Validar que la respuesta contenga status 200.\\
		4. Consultar recurso para confirmar cambios.\\
		5. Verificar que updated\_at se modifica.
		}
		\testcaseexpresult{
		- Status HTTP: 200 OK.\\
		- Metadatos actualizados en BD.\\
		- Historial versionado (opcional).\\
		- Cambios persistidos.\\
		- Campo updated\_at actualizado.
		}
		\testcaseuserstory{14}
		\end{acceptancetest}

		\begin{acceptancetest}[tb:HU14_P7]
		\testcasecode{HU14\_P7}
		\testcasename{Listar recursos de feature con filtros}
		\testcasedescription{Validar que se pueden obtener todos los recursos de una feature con filtros.}
		\testcaseexeccond{
		- Feature tiene múltiples recursos asociados.\\
		- Usuario está autenticado.
		}
		\testcaseexecstep{
		1. Enviar GET a /api/v1/models/\{model\_id\}/versions/\{version\_id\}/features/\{feature\_id\}/resources?language=es\&type=VIDEO con headers.\\
		2. Validar que la respuesta contenga status 200.\\
		3. Verificar que solo recursos con idioma es y tipo VIDEO se retornan.\\
		4. Cada recurso incluye metadatos completos.\\
		5. Paginación funciona.
		}
		\testcaseexpresult{
		- Status HTTP: 200 OK.\\
		- Array de recursos filtrados.\\
		- Información completa de cada uno.\\
		- Filtros aplicados correctamente.\\
		- Paginación funciona.
		}
		\testcaseuserstory{14}
		\end{acceptancetest}

\end{addendum}

	\addappendix[bloque]{Anexos/BloqueCoronaReport}
%\addappendix[controlador]{Anexos/ControladorIndustrialReport}
